\documentclass[12pt,a4paper]{article}
\usepackage[utf8]{inputenc}
\usepackage[T1]{fontenc}
\usepackage{amsmath,amssymb,amsfonts}
\usepackage{graphicx}
\usepackage{booktabs}
\usepackage{xcolor}
\usepackage[
    colorlinks=true,
    linkcolor=blue!70!black,
    citecolor=green!50!black,
    urlcolor=blue!70!black,
    bookmarks=true,
    bookmarksopen=true,
    pdfauthor={Razvan-Constantin Anghelina},
    pdftitle={Geometric Foundations of Universal Constants: A 600-Cell Framework for Coupling Constants, Mass Hierarchies, and Gauge Unification},
    pdfsubject={Theoretical Physics},
    pdfkeywords={600-cell, fine structure constant, gauge unification, golden ratio, Planck scale}
]{hyperref}
\usepackage{geometry}
\usepackage{float}
\usepackage{caption}
\usepackage{subcaption}
\usepackage{multirow}

\geometry{margin=2.5cm}

\title{\textbf{Why Three Generations?\\
A Derivation from the 600-Cell Polytope}}

\author{
Razvan-Constantin Anghelina\\
\texttt{contact@webdesignmedia.ro}\\[0.3em]
Interactive visualization: \url{https://webdesignmedia.ro/hjer_particles/}\\[1em]
\textit{Version 4.1 -- February 2026}
% v4.1: McKay correspondence, PMNS predictions, testable predictions tiers
% v4.0: Pyramid restructure (strongest results first)
% v3.4: Edge split 1+3+8, Gram=24-cell, A_eff=A^2/3 (EXP-143,144)
}

\date{}

\begin{document}

\maketitle

\begin{abstract}
We present a framework in which the fundamental constants of particle physics emerge from the geometric and spectral properties of the 600-cell, the unique universal energy minimizer in $\mathbb{R}^4$ (Cohn--Kumar theorem). Its spectral properties encode coupling constants with high precision: $\alpha$ (0.0001\%), $\alpha_s$ (0.03\%), $\sin^2\theta_W$ (0.19\%), and $m_H$ (0.09\%). The 24 non-fermionic vertices form the $D_4$ root system, yielding $SU(3) \times SU(2) \times U(1)$ with a natural $1 + 3 + 8 = 12$ edge splitting that matches $\dim(\text{SM})$. The remaining 96 vertices encode fermions ($96 = 16 \times 3 \times 2$). A unified mass formula $m_f = m_e \cdot \phi^{a \cdot d_{\text{eff}} + b \cdot k_{\text{eff}}}$ predicts 7/9 fermion masses within 2\%, where $(a,b)$ are uniquely fixed by complexity minimization on the graph.

\textbf{Central result:} We prove that the number of fermion generations is \textbf{exactly three}. The proof proceeds in five steps: (1)~the 600-cell spectrum places mass levels in $\mathbb{Z}[\phi]$, with $z = a + b\phi$; (2)~the $E_8$ icosian construction gives a Galois shadow $z' = \sigma(z)$ in the conjugate sector; (3)~since $E_8$ is defined over $\mathbb{Z}$, the energy functional must be Galois-invariant, forcing $E = V(z) + V(z')$, whose vanishing requires $z$ to be a unit of $\mathbb{Z}[\phi]$; (4)~by Dirichlet's unit theorem, units on the $a=1$ line satisfy $F(k\!-\!1) = 1$, which has exactly three Fibonacci solutions; (5)~therefore $b \in \{0,1,2\}$, giving three generations. No free parameters enter this derivation.

The $E_8$ root system is constructed explicitly via the icosian lattice ($S \cup T$, $T = \phi' S$, 240 roots exact). The McKay correspondence maps the 9 irreps of the binary icosahedral group $2I$ to the affine $E_8$ Dynkin diagram, and the CKM bare exponents $\{3,7,12\}$ are expressed as $n = 9b_2 - 5b_1 - 6$, where the coefficients $(9, 5, 6)$ equal the dimensions of the three $E_8$ Dynkin legs. The edge structure forbids gauge--gauge couplings, exhibits $U(1)$ topological confinement (perfect matching, gap $= 2$), and supports pure $SU(3)$ plaquettes reproducing the non-Abelian self-interaction of QCD. Topological solitons on the graph have uniform energy $E_{\text{flip}} = 3$, and Kibble--Zurek cooling freezes ${\sim}88$ defects---a possible cosmological matter origin. All three CKM mixing angles are predicted to $< 0.2\%$: $\theta_{12} = \arctan(\phi^{-3}(1 - 2\alpha_s/3\pi))$ to $0.001\%$, $\theta_{23}$ to $0.17\%$, and $\theta_{13}$ to $0.096\%$, combining bare geometric values ($\phi^{-n}$ with $n \in \{3,7,12\}$) with perturbative gauge corrections. All three PMNS mixing angles are predicted within $1\sigma$ of PDG~2024 values using purely geometric formulas: $\sin^2\theta_{13} = 1/45 = 1/(\text{diam} \times N_{\text{eig}})$ to $0.16\%$, $\sin^2\theta_{12} = 2/(\phi+5)$ to $0.33\%$, and $\sin^2\theta_{23} = 4/7$ to $0.32\%$. The mass ratio corrections exhibit Galois complementarity: lepton corrections vanish in the $\phi$-sector and live entirely in the Galois shadow ($\phi'$-sector), while up-quark corrections show the inverse pattern with $\delta_k \approx \alpha_s$ to $0.5\%$.

\textbf{Key testable prediction:} $\sin^2\theta_{13}(\text{PMNS}) = 1/45 = 0.02222$, predicted to $0.16\%$ precision. The JUNO experiment will measure this quantity with sub-percent accuracy, providing a direct test of the geometric origin.
\end{abstract}

\tableofcontents
\newpage

%==============================================================
\section{Introduction}
%==============================================================

The search for a fundamental structure underlying the constants of nature has been a central pursuit in theoretical physics. The Standard Model of particle physics, while remarkably successful, contains approximately 19 free parameters whose values must be determined experimentally. The origin of these numbers---why the fine structure constant equals $1/137.036$ rather than some other value---remains one of the deepest mysteries in physics.

We propose that spacetime at the Planck scale possesses the discrete structure of the 600-cell (hexacosichoron), a regular four-dimensional polytope with exceptional geometric properties. This approach is motivated by the Cohn-Kumar theorem \cite{cohn2007,cohn2011}, which establishes the 600-cell as the unique universal energy minimizer in $\mathbb{R}^4$ for a broad class of potentials. The geometric framework builds upon the 600-cell GUT developed by O'Neill \cite{oneill2025gut,oneill2023binary}.

The key insight is that the fundamental constants are not arbitrary but emerge as ``spectral fingerprints'' of this optimal geometric structure. The golden ratio $\phi = (1+\sqrt{5})/2$, which permeates the 600-cell geometry, appears naturally in the coupling constants and mixing angles of particle physics.

This paper is organized as follows: Section~2 describes the 600-cell architecture and its spectral properties. Section~3 derives the Standard Model gauge group from the $D_4$ root system, including interaction rules and the $1+3+8$ edge splitting. Section~4 derives coupling constants and the Higgs mass. Section~5 maps fermions to the vertex structure, \emph{proves} that the number of generations is exactly three via Galois invariance and the Fibonacci condition, and derives the democratic bare mixing matrix from representation branching. Section~6 presents the unified fermion mass formula. Section~7 identifies Discrete Scale Invariance as the underlying mechanism. Section~8 constructs the $E_8$ root system from the icosian lattice and establishes the McKay correspondence, including the CKM exponent formula. Section~9 validates the framework and proposes tiered experimental tests. Section~10 summarizes results with an explicit discussion of limitations. Appendix~\ref{app:ckm} presents the mixing matrices, including new PMNS predictions.

%==============================================================
\section{The 600-Cell Architecture}
%==============================================================

\subsection{Fundamental Properties}

The 600-cell is a regular convex 4-polytope with the following characteristics:

\begin{table}[H]
\centering
\begin{tabular}{lrl}
\toprule
\textbf{Property} & \textbf{Value} & \textbf{Physical Interpretation} \\
\midrule
Vertices & 120 & Lattice nodes / degrees of freedom \\
Edges & 720 & Propagation paths \\
Faces & 1200 & Triangular interfaces \\
Cells & 600 & Tetrahedral units \\
Symmetry group & $H_4$ & Order 14,400 \\
Vertex degree & 12 & Neighbors per node \\
Graph diameter & 5 & Maximum geodesic distance \\
Euler characteristic & 0 & Topological invariant \\
\bottomrule
\end{tabular}
\caption{Properties of the 600-cell polytope.}
\end{table}

\subsection{Universal Optimality}

\textbf{Theorem (Cohn-Kumar \cite{cohn2007}):} The 600-cell is the unique configuration in $\mathbb{R}^4$ that acts as a universal energy minimizer for completely monotonic potentials. The analogous result for $E_8$ in $\mathbb{R}^8$ was established by Cohn, Kumar, Miller, Radchenko, and Viazovska \cite{cohn2022}, strengthening the connection between the 600-cell and $E_8$ (Section~\ref{sec:e8}).

\textbf{Physical consequence:} By the principle of least action, spacetime at the Planck scale self-organizes into the 600-cell lattice structure, analogous to:
\begin{itemize}
    \item 2D: Hexagonal lattice (honeycomb)
    \item 3D: FCC/HCP crystal structures
    \item 4D: 600-cell
\end{itemize}

Elser and Sloane \cite{elser1987} showed that projecting the $E_8$ lattice onto 4D via icosahedral symmetry produces a quasicrystal built from two concentric 600-cells in the ratio $\phi$---the same structure that emerges from the icosian lattice construction (Section~\ref{sec:e8}).

\subsection{Vertex Decomposition}

The 120 vertices decompose into three geometrically distinct sets:

\begin{equation}
120 = 8 + 16 + 96
\end{equation}

\begin{table}[H]
\centering
\begin{tabular}{ccll}
\toprule
\textbf{Set} & \textbf{Count} & \textbf{Structure} & \textbf{Physical Role} \\
\midrule
Type A & 8 & 16-cell (cross-polytope) & \multirow{2}{*}{$D_4$ root system $\to$ gauge sector (Section~\ref{sec:gauge})} \\
Type B & 16 & 8-cell (tesseract) & \\
Type C & 96 & Snub 24-cell & Fermions ($16 \times 3 \times 2$) \\
\bottomrule
\end{tabular}
\caption{Vertex decomposition and physical interpretation.}
\end{table}

The 96 golden-ratio vertices encode:
\begin{equation}
96 = 16 \times 3 \times 2 = \text{(fermions per generation)} \times \text{(generations)} \times \text{(chiralities)}
\end{equation}

\subsection{Spectral Structure and Association Scheme}
\label{sec:spectral}

The 600-cell graph is \textbf{distance-transitive}: vertices at the same graph distance are related by symmetry. This creates an \textbf{association scheme} with 9 angular distance classes $d \in \{0, 1, \ldots, 8\}$ (corresponding to $S^3$ arc distances $\theta = 0, \pi/5, \ldots, \pi$), and the adjacency matrix has exactly 9 distinct eigenvalues---all of the form $a + b\phi$ with $a,b \in \mathbb{Z}$:

\begin{table}[H]
\centering
\begin{tabular}{cccl}
\toprule
\textbf{Index} $i$ & \textbf{Eigenvalue} $\theta_i$ & \textbf{Multiplicity} $m_i$ & \textbf{Structure} \\
\midrule
0 & 12             & $1 = 1^2$   & Trivial representation \\
1 & $6\phi$        & $4 = 2^2$   & \multirow{2}{*}{Galois pair in $\mathbb{Q}(\sqrt{5})$} \\
8 & $6 - 6\phi$    & $4 = 2^2$   & \\
2 & $4\phi$        & $9 = 3^2$   & \multirow{2}{*}{Galois pair} \\
6 & $4 - 4\phi$    & $9 = 3^2$   & \\
3 & 3              & $16 = 4^2$  & \multirow{2}{*}{Rational pair} \\
7 & $-3$           & $16 = 4^2$  & \\
4 & 0              & $25 = 5^2$  & Kernel \\
5 & $-2$           & $36 = 6^2$  & Largest eigenspace \\
\bottomrule
\end{tabular}
\caption{Adjacency spectrum of the 600-cell. All multiplicities are perfect squares $n^2$ for $n = 1, \ldots, 6$, reflecting $H_4$ irreducible representations.}
\label{tab:spectrum}
\end{table}

The characteristic polynomial over $\mathbb{Q}$ factors as:
\begin{equation}
P(x) = (x-12)(x^2-6x-36)^4(x^2-4x-16)^9(x-3)^{16}\,x^{25}(x+2)^{36}(x+3)^{16}
\end{equation}

Two key structural features:
\begin{enumerate}
    \item \textbf{Intersection numbers:} Each edge of the 600-cell has exactly $a_1 = 5$ common neighbors and $b_1 = 6$ next-shell neighbors. These are topological invariants---identical for all 720 edges.
    \item \textbf{Phi-sector:} The four eigenspaces with $\phi$-dependent eigenvalues ($i = 1, 2, 6, 8$) have total dimension $4 + 9 + 9 + 4 = 26$. This is the same 26 that appears in $\sin^2\theta_W = 6/26$ (Section~4.3).
\end{enumerate}

Three exact \textbf{Laplacian eigenvalue ratios} connect the spectrum to physics:
\begin{align}
\lambda_4 / \lambda_1 &= 2\phi^2 &\quad &\text{(yields } \alpha \text{, Section~4.1)} \\
\lambda_6 / \lambda_2 &= \phi^2 &\quad &\text{(connects to QED vertex corrections)} \\
\lambda_7 / \lambda_3 &= 5/3 = F/E &\quad &\text{(faces/edges = plaquettes per link, EXP-132)}
\end{align}

where $\lambda_i = 12 - \theta_i$ are Laplacian eigenvalues.


%==============================================================
\section{Gauge Group from the $D_4$ Root System}
%==============================================================

\subsection{The Gauge Sector: $D_4$ Root System from the 24-Cell}
\label{sec:gauge}

The 24 non-fermionic vertices (Type~A + Type~B) form a \textbf{24-cell}, which is simultaneously a regular 4-polytope and a root system. This root system is $D_4$.

\textbf{Verification.} The 8 Type-A vertices $(\pm 1, 0, 0, 0)$, etc., together with the 16 Type-B vertices $(\pm\tfrac{1}{2})^4$, satisfy the Cartan integer condition $2\langle \alpha, \beta\rangle / \langle \alpha, \alpha\rangle \in \mathbb{Z}$ for all root pairs, confirming they form the $D_4$ root system (rank 4, 24 roots).

\textbf{Standard Model gauge group.} The maximal-rank regular subalgebra decomposition of $D_4$ gives:
\begin{equation}
\boxed{D_4 \;\to\; A_2 \oplus A_1 \oplus U(1) \;=\; SU(3) \times SU(2) \times U(1)}
\end{equation}

This is a standard result in Lie theory. The dimension and rank match exactly:
\begin{align}
\dim(SU(3) \times SU(2) \times U(1)) &= 8 + 3 + 1 = 12 = \text{vertex degree of 600-cell} \\
\mathrm{rank}(SU(3) \times SU(2) \times U(1)) &= 2 + 1 + 1 = 4 = \dim(\mathbb{R}^4)
\end{align}

The remaining 96 vertices (Type~C, the snub 24-cell) naturally separate as the \textbf{matter sector}, with $96 = 16 \times 3 \times 2$ encoding 16 fermion states per generation, 3 generations, and 2 chiralities.

\textbf{Group-theoretic origin of the gauge/matter split (EXP-154).} The partition $24 + 96$ has a rigorous origin in the \emph{binary icosahedral group} $2I$ (order~120), which acts on the 600-cell by left quaternion multiplication---in fact, the 600-cell graph \emph{is} the Cayley graph of $2I$ with respect to its 12 generators of order~10. The 24 gauge vertices form the \emph{binary tetrahedral subgroup} $2T \subset 2I$ (order~24), verified to be closed under quaternion multiplication. The five left cosets of $2T$ in $2I$ yield five inscribed 24-cells:
\begin{equation}
\text{Coset}_0 = 2T = \text{Type A+B} \;(8+16), \qquad \text{Cosets}_{1\text{--}4} = 4 \times 24 \;\text{Type-C vertices}
\end{equation}
Each coset forms a valid 24-cell (degree~8, 96~edges). The edge structure is \textbf{completely 5-partite}: zero edges connect vertices within the same coset, and exactly 72 edges connect each pair of distinct cosets ($\binom{5}{2} \times 72 = 720$). Left multiplication by $2I$ permutes the 5~cosets, yielding a homomorphism $2I \to A_5$ (kernel~$= Z_2$, image~$= A_5$ since all coset permutations are even). The 6 decagons through each vertex correspond to the 6~Sylow 5-subgroups of the vertex-figure symmetry group~$A_5$.

\textbf{$D_4$ triality.} The Dynkin diagram of $D_4$ has an $S_3$ outer automorphism group (triality), connecting three inequivalent 8-dimensional representations: $\mathbf{8}_v$, $\mathbf{8}_s$, $\mathbf{8}_c$. Left quaternion multiplication by $\omega = \tfrac{1}{2}(1,1,1,1)$ is a full 600-cell automorphism that cyclically permutes $\mathbf{8}_v \to \mathbf{8}_s \to \mathbf{8}_c \to \mathbf{8}_v$, while preserving the 96 Type-C (matter) vertices. Consequently, the apparent $8 + 16$ split (Type~A versus Type~B) is a \emph{coordinate artifact}: the triality symmetry is \textbf{fully preserved} by the 600-cell geometry (EXP-126).

\textbf{$A_5$ irrep selection.} Although the 600-cell alone cannot break triality, the \emph{vertex figure} (icosahedron, symmetry $A_5$) provides a partial selection. The 12 edges at each vertex decompose under $A_5$ as $12 = 1 + 3 + 3' + 5$. The \emph{unique} regrouping compatible with SM dimensions is $8 = 3' + 5$, $3 = 3$, $1 = 1$, corresponding to $SU(3) \times SU(2) \times U(1)$. The competing decomposition $SU(2)^4$ would require $3 + 3 + 3 + 3$, which is incompatible with the $A_5$ structure.

\textbf{$E_8$ triality breaking (EXP-130, 133).} The $E_8$ embedding (Section~\ref{sec:e8}) resolves the triality question definitively. The triality automorphism $\omega$ maps all 240 $E_8$ roots to $E_8$ roots and preserves both 600-cell copies $S$ and $T$ separately---so at the \emph{root system} level, triality is an $E_8$ symmetry. However, at the \emph{Dynkin diagram} level, triality is \textbf{broken}. The $D_4$ sub-diagram $\{2, 3, [4], 5\}$ sits inside the $E_8$ Dynkin diagram (Bourbaki numbering):
\begin{equation}
1 - 3 - \boxed{4} - 5 - 6 - 7 - 8 \qquad \text{with} \quad 2 - \boxed{4}
\end{equation}
where node $[4]$ is the $D_4$ branching node. The three legs extend into $E_8$ with \textbf{unequal tail lengths}, and the Kac labels (marks of the highest root) are all \textbf{distinct}:

\begin{table}[H]
\centering
\begin{tabular}{ccccl}
\toprule
\textbf{D4 Leg} & \textbf{Tail length} & \textbf{Kac label} & \textbf{Tail algebra} & \textbf{SM assignment} \\
\midrule
Node 5 $\to$ 6,7,8 & 3 & $m = 5$ & $A_4 \supset A_2$ & $SU(3)_{\text{color}}$ \\
Node 3 $\to$ 1       & 1 & $m = 4$ & $A_1$              & $SU(2)_{\text{weak}}$ \\
Node 2 (dead end)     & 0 & $m = 3$ & ---                & $U(1)_Y$ (via $SU(2)_R$) \\
\bottomrule
\end{tabular}
\caption{$D_4$ leg assignment from $E_8$ Dynkin diagram topology. Tail length: number of $E_8$ nodes beyond $D_4$ on each leg.}
\end{table}

The assignment is \textbf{unique and monotonic}: longer tail $\to$ larger gauge group $\to$ more generators. Since $E_8$ has no nontrivial outer automorphisms, the triality can only act as an inner automorphism (Weyl group element), which assigns distinct roles to the three legs via the Weyl chamber structure. The standard GUT breaking chain is fully recovered:
\begin{equation}
E_8 \;\to\; E_6 \times SU(3)_{\text{fam}} \;\to\; SO(10) \;\to\; SU(5) \;\to\; SU(3)_C \times SU(2)_L \times U(1)_Y
\end{equation}

The $D_4$ simple roots within $E_8$ come from \emph{mixed} 600-cells: 2 from $S_{96}$ and 2 from $T_{96}$ (snub 24-cell vertices of each icosian copy), confirming that the two 600-cells play asymmetric roles in the $E_8$ root structure.

\textbf{$SU(2)_R \to U(1)_Y$ breaking (EXP-135).} The dead-end leg (node 2) generates $SU(2)_R$, which reduces to $U(1)_Y$ through three independent structural mechanisms:
\begin{enumerate}
    \item \textbf{Topological isolation:} The dead end has no non-Abelian extension in $E_8$ to protect $SU(2)_R$ from breaking.
    \item \textbf{Maximal commutant:} Removing node 2 from $E_8$ leaves $A_7 = SU(8)$ (dim~63)---the largest surviving sub-algebra among the three D4 legs.
    \item \textbf{Smallest Kac label:} $m_2 = 3$ is the minimum among D4-leg Kac labels, giving the most constrained $U(1)$ charge spectrum ($q_Y \in \{0, \pm 1, \pm 2, \pm 3\}$).
\end{enumerate}
The $U(1)_Y$ charge spectrum of the 120 positive $E_8$ roots decomposes as $28 + 56 + 28 + 8$ for charges $q_Y = 0, 1, 2, 3$ respectively.

\textbf{Classification:} The full chain 600-cell $\to$ 24-cell $\to$ $D_4$ $\to$ $E_8$ Dynkin $\to$ $SU(3) \times SU(2) \times U(1)$ is \textbf{DERIVED} from geometric and algebraic structure, with no free parameters. The identification of the 96-vertex complement as the matter sector is a PATTERN (the correct counting but no spinor representation derivation).

\subsection{Interaction Selection Rules and Propagator (EXP-136, 137)}
\label{sec:interactions}

The edge structure of the 600-cell imposes \textbf{automatic selection rules} on particle interactions. Since no edges connect gauge vertices to each other:
\begin{equation}
\text{A-A} = \text{A-B} = \text{B-B} = 0 \quad \text{(edges)}
\end{equation}
the only non-vanishing edges are A-C (96), B-C (192), and C-C (432), where C denotes the 96 fermion vertices. This propagates to interaction vertices: of the 10 possible triangle types and 15 possible tetrahedron types, only \textbf{three} survive in each case:

\begin{table}[H]
\centering
\begin{tabular}{lcccl}
\toprule
\textbf{Triangles} & \textbf{Count} & \textbf{Per fermion} & \textbf{Ratio} & \textbf{Physical analogue} \\
\midrule
ACC (gauge--fermion--fermion) & 240 & 5 & 1 & QED/QCD vertex \\
BCC (Higgs--fermion--fermion) & 480 & 10 & 2 & Yukawa coupling \\
CCC (fermion--fermion--fermion) & 480 & 15 & 3 & Fermion self-coupling \\
\midrule
\textbf{Total} & \textbf{1200} & \textbf{30} & & \\
\bottomrule
\end{tabular}
\caption{Triangle (3-point vertex) classification. Per-fermion counts are uniform across all 96 Type-C vertices. The 1:2:3 ratio is exact.}
\end{table}

Similarly, tetrahedra (4-point vertices) decompose as ACCC\,(160) : BCCC\,(320) : CCCC\,(120), with ratio 2:4:1.5 reflecting the same hierarchy.

\textbf{Key consequence:} Gauge bosons interact \emph{only} through fermion mediators. The ``effective gauge adjacency'' (2-step paths through fermions) connects every gauge pair by exactly 3 paths, with eigenvalues $\{24, 12, 0, -6, -12\}$---all integers despite the irrational ($\phi$-dependent) spectrum of the full graph.

\textbf{Propagator structure.} The massless Green's function $G = (L + 0)^{-1}_{\perp}$ (pseudoinverse of the Laplacian, projecting out the constant mode) shows \textbf{oscillatory behavior}:

\begin{center}
\begin{tabular}{crl}
$d$ & $\langle G_{\text{CC}}(d)\rangle$ & Sign \\
\hline
1 & $+0.01527$ & $+$ \\
2 & $+0.00236$ & $+$ \\
3 & $-0.00335$ & $-$ \\
4 & $-0.00634$ & $-$ \\
5 & $-0.00791$ & $-$ \\
\end{tabular}
\end{center}

The sign change at $d = 3$ is characteristic of \textbf{confinement} on finite graphs: beyond the interaction range of common gauge neighbors (which drops to zero at $d \geq 3$), the propagator becomes repulsive. At $d = 1$, the first eigenspace (Laplacian eigenvalue $\lambda_1 = 12 - 6\phi \approx 2.29$) contributes 77\% of the total amplitude, confirming that the lowest spectral mode dominates short-range interactions. The propagator values live in $\mathbb{Q}(\sqrt{5})$; for example, $G(1) = (-4 + 23\phi)/2175$.

\textbf{Classification:} Selection rules and triangle/tetrahedron counts are \textbf{DERIVED}. The mapping to SM interaction types (ACC $\to$ gauge vertex, BCC $\to$ Yukawa, CCC $\to$ self-coupling) is \textbf{PATTERN}. The confinement interpretation is \textbf{SPECULATIV}.

\subsection{The $1 + 3 + 8$ Edge Split and Effective 24-Cell (EXP-143, 144)}
\label{sec:edge_split}

Each fermion vertex has 12 neighbors: 1 of type A, 2 of type B, and 9 of type C. The Standard Model gauge group has dimension $1 + 3 + 8 = 12$, matching the total but not the internal structure ($1 + 2 + 9$ vs.\ $1 + 3 + 8$). A detailed analysis of the local neighborhood reveals a \textbf{natural geometric splitting} that resolves this discrepancy.

\textbf{Unique special CC edge.} Among the 9 C-neighbors of each fermion, exactly \textbf{one} has a distinct gauge-overlap profile:
\begin{itemize}
    \item Degree 4 in the local CC subgraph (vs.\ degree 3 for 6 others)
    \item Shares the same A-vertex (cross-polytope) gauge neighbor as the central fermion
    \item Shares \textbf{zero} B-vertex (tesseract) gauge neighbors
    \item Has minimum total gauge overlap (1 vs.\ 2 for the others)
\end{itemize}
This classification is \textbf{uniform} across all 96 fermion vertices, with exactly 48 such ``special'' CC edges in total. The per-fermion edge profile is:

\begin{center}
\begin{tabular}{ccccc}
\toprule
\textbf{Edge type} & \textbf{Count} & \textbf{shared A} & \textbf{shared B} & \textbf{SM sector} \\
\midrule
AC & 1 & -- & -- & $U(1)_Y$ \\
BC & 2 & -- & -- & $SU(2)_L$ (charged) \\
CC$_{\text{special}}$ & 1 & 1 & 0 & $SU(2)_L$ (neutral) \\
CC$_{\text{regular}}$ & 8 & mixed & mixed & $SU(3)_C$ \\
\midrule
\textbf{Total} & \textbf{12} & & & $\mathbf{1 + 3 + 8}$ \\
\bottomrule
\end{tabular}
\end{center}

\textbf{Effective gauge adjacency = 24-cell.} The Gram matrix $G = M_{\text{GF}} M_{\text{GF}}^T$, where $M_{\text{GF}}$ is the $24 \times 96$ gauge--fermion incidence matrix, encodes the effective gauge--gauge coupling mediated by fermions. Its off-diagonal entries take only values 0 or 3 (shared fermion neighbors), so $G = 12\mathbf{I} + 3\,A_{\text{eff}}$ where $A_{\text{eff}}$ has entries in $\{0, 1\}$. Remarkably:
\begin{equation}
A_{\text{eff}} = A_{\text{24-cell}} \quad \text{(verified element-by-element)}
\end{equation}
The effective adjacency through shared fermions is \textbf{identical} to the nearest-neighbor adjacency of the standard 24-cell (D$_4$ root system). Furthermore, this arises from the full 600-cell adjacency $A$ via:
\begin{equation}
A_{\text{eff}} = \frac{1}{3}\, A^2\big|_{\text{gauge}} \,,
\end{equation}
where $A^2|_{\text{gauge}}$ restricts the square of the adjacency matrix to the 24 gauge vertices. The factor $1/3$ is exact and universal.

\textbf{Hop-invariant eigenspaces.} The Gram matrix spectrum is $\{36, 24, 12, 6, 0\}$ with multiplicities $\{1, 4, 9, 8, 2\}$---all eigenvalues are multiples of 6, and the first three multiplicities are perfect squares. Higher-hop matrices $G_n = M_{\text{GF}} A_{\text{CC}}^n M_{\text{GF}}^T$ share the \textbf{same eigenspaces} for all $n$, with the eigenvalue $6$ (multiplicity 8) remaining \textbf{hop-invariant}. The dominant eigenvalue follows $36 \cdot 9^n$, where 9 is the CC-subgraph degree.

\textbf{Plaquette structure (EXP-145).} The $1 + 3 + 8$ edge classification induces a decomposition of all 1200 triangular plaquettes into five types:
\begin{center}
\begin{tabular}{lcc}
\toprule
\textbf{Type} & \textbf{Count} & \textbf{Pure/Mixed} \\
\midrule
$SU(3)$-$SU(3)$-$SU(3)$ & 288 & Pure \\
$SU(2)$-$SU(2)$-$SU(3)$ & 480 & Mixed \\
$SU(2)$-$SU(3)$-$SU(3)$ & 192 & Mixed \\
$SU(3)$-$U(1)$-$U(1)$ & 192 & Mixed \\
$SU(2)$-$U(1)$-$U(1)$ & 48 & Mixed \\
\bottomrule
\end{tabular}
\end{center}
Only $SU(3)$ has pure plaquettes; $U(1)$ and $SU(2)$ have \textbf{zero}. This is the lattice analogue of the fact that only gluons carry color charge and exhibit cubic self-interaction---the photon and $W/Z$ bosons do not self-interact in the pure gauge sector. The 288 pure $SU(3)$ plaquettes satisfy $288 = 24 \times 12 = \operatorname{Tr}(G)$.

\textbf{Vertex-transitivity and coupling independence (EXP-146).} The $H_4$ symmetry group acts transitively on the 120 vertices, with a fundamental consequence for the spectral action: every eigenspace of the adjacency matrix projects \textbf{identically} onto all three gauge sectors:
\[
\chi_k(\text{U1}) = \chi_k(\text{SU2}) = \chi_k(\text{SU3}) \quad \forall\, k = 0, \ldots, 8,
\]
where $\chi_k(S) = |E_S|^{-1} \sum_{(i,j) \in E_S} P_k[i,j]$ is the average projection of eigenspace $k$ onto edges of sector $S$. Similarly, the vertex-type weights are proportional to $|A|:|B|:|C| = 8:16:96$ in every eigenspace without exception. Consequently, no Connes--Chamseddine spectral action $\operatorname{Tr}(f(D^2))$ can differentiate coupling constants across sectors---the edge-type structure is invisible to the spectrum.

This establishes a \textbf{separation theorem}: the 600-cell encodes two independent geometric properties that correspond to two independent physical structures:
\begin{enumerate}
\item \textbf{Lagrangian form} (from topology): the $1 + 3 + 8$ edge split, plaquette types, selection rules---determining \emph{what} interacts with \emph{what}.
\item \textbf{Coupling values} (from the spectrum): eigenvalue ratios $\theta_1/\theta_3 = 2\phi$ and the intersection number $a_1 = 5$---determining \emph{how strongly} they interact.
\end{enumerate}
This mirrors the Standard Model, where the gauge group structure and coupling constant values are also logically independent.

\textbf{Shared-neighbor asymmetry (EXP-147).} While the spectral projection is sector-blind, the \emph{topological} neighborhood structure distinguishes sectors. Each edge has exactly $a_1 = 5$ shared neighbors (distance-regular), but their type decomposition is \textbf{exact} (standard deviation zero within each edge type):
\begin{center}
\begin{tabular}{lcccc}
\toprule
\textbf{Edge type} & \textbf{Count} & \textbf{shared A} & \textbf{shared B} & \textbf{shared C} \\
\midrule
$U(1)$ (AC) & 96 & 0 & 0 & 5 \\
$SU(2)_{\text{ch}}$ (BC) & 192 & 0 & 0 & 5 \\
$SU(2)_{\text{n}}$ (CC$_*$) & 48 & 1 & 0 & 4 \\
$SU(3)$ sub-type 1 & 96 & 0 & 1 & 4 \\
$SU(3)$ sub-type 2 & 96 & 0 & 2 & 3 \\
$SU(3)$ sub-type 3 & 192 & 1 & 1 & 3 \\
\bottomrule
\end{tabular}
\end{center}
$U(1)$ and $SU(2)_{\text{ch}}$ edges share \textbf{only fermion} neighbors (zero gauge-boson access), while $SU(3)$ edges share 1.75 gauge-boson neighbors on average across three sub-types. This is the lattice realization of Abelian vs.\ non-Abelian character: the photon interacts only through fermion loops, while gluons have direct gauge-boson access.

\textbf{Per-fermion gauge budget.} Each fermion has a total gauge access of exactly 15 (uniform across all 96 fermions): 0 from $U(1)$ edges, 1 from $SU(2)$ edges (solely from the neutral CC$_*$), and 14 from $SU(3)$ edges.

\textbf{Topological decomposition of $10\phi$.} The total fermion-shared neighbors across all $U(1)$ edges is $96 \times 5 = 480$, while the total gauge-shared neighbors across all $SU(2)$ edges is $48 \times 1 = 48$ (only the neutral CC$_*$ edges contribute). Their ratio is:
\begin{equation}
\frac{\text{fermion-shared}_{U(1)}}{\text{gauge-shared}_{SU(2)}} = \frac{480}{48} = 10\,.
\end{equation}
Since $\alpha_s/\alpha = 10\phi$ (Section~\ref{sec:alpha_s}), the factor 10 now has a topological origin: it counts how many fermion-loop mediations in the $U(1)$ sector correspond to one gauge self-interaction in the $SU(2)$ sector. The factor $\phi$ remains spectral (from $\theta_1/\theta_3 = 2\phi$, where the factor 2 itself equals $n_{U(1)}/n_{\text{CC}_*} = 96/48$). This provides a dual interpretation of the strong-to-electromagnetic coupling ratio: spectral ($a_1 \cdot \theta_1/\theta_3$) and topological ($\text{fermion}_{U(1)}/\text{gauge}_{SU(2)} \cdot \phi$).

\textbf{$SU(3)$ sub-type decomposition (EXP-148).} The 384 $SU(3)$ edges decompose into three sub-types based on their shared-neighbor profile, each forming a distinct sub-graph on the 96 fermion vertices:
\begin{center}
\begin{tabular}{lccccc}
\toprule
\textbf{Sub-type} & \textbf{Edges} & \textbf{Degree} & \textbf{Components} & \textbf{Per fermion} & \textbf{CC$_*$ link} \\
\midrule
T1 (0A, 1B, 4C) & 96 & 2 & 24 & 2 & 100\% \\
T2 (0A, 2B, 3C) & 96 & 2 & 32 & 2 & 0\% \\
T3 (1A, 1B, 3C) & 192 & 4 & \textbf{8} & 4 & 50\% \\
\bottomrule
\end{tabular}
\end{center}
The sub-type T3 has exactly \textbf{8 connected components}, each containing 12 fermion vertices---matching $\dim(SU(3)) = 8$ and the vertex degree 12. The mechanism is \emph{derived}: each T3 edge shares exactly one Type-A neighbor (by its profile $(1A,1B,3C)$), and the 8 Type-A vertices partition the 192 T3 edges into 8 disjoint groups of 24, each forming a single connected component (EXP-149). Thus $\dim(SU(3)) = |A| = 8$ is a geometric identity, not a numerical coincidence. Similarly, T1 components are labeled by the 24 A--B pairs each edge accesses, and T2 components by the 32 B--B pairs. The three sub-type adjacency matrices do \textbf{not commute} ($[A_{T_i}, A_{T_j}] \neq 0$ for all pairs), encoding non-Abelian structure consistent with $SU(3)$.

The CC$_*$ connectivity reveals a hierarchy: T1 edges are \emph{all} connected to the $SU(2)$ neutral partner (linking the strong and weak sectors), T2 edges are \emph{completely disconnected} from it, and T3 edges are connected at 50\%. The 288 pure $SU(3)$ plaquettes decompose as $192\,(\text{T1-T2-T3}) + 64\,(\text{T3-T3-T3}) + 32\,(\text{T2-T2-T2})$; the dominant mixed type reflects inter-sub-type coupling.

\textbf{Classification:} The $1 + 3 + 8$ edge split, effective adjacency $A_{\text{eff}} = A_{\text{24-cell}}$, plaquette decomposition, separation theorem, shared-neighbor profiles, the topological decomposition of $10\phi$, and the $SU(3)$ sub-type structure including the A-vertex generation of 8 components are \textbf{DERIVED}. The SM sector assignments remain \textbf{PATTERN}.


\subsection{U(1) Confinement from Perfect Matching (EXP-167)}
\label{sec:u1_confinement}

The 48 $U(1)$ edges (CC$_*$ type, shared-neighbor profile $(1A, 0B)$) exhibit an extreme structural property: they form a \textbf{perfect matching} on the 96 fermion vertices.

\textbf{Definition.} A perfect matching is a set of edges in which every vertex appears exactly once. Equivalently, the $U(1)$ subgraph has degree 1 at every fermion vertex.

\textbf{Consequences:}
\begin{itemize}
    \item The $U(1)$ subgraph consists of 48 isolated edges (48 connected components, each of size 2).
    \item The cycle rank is zero: the $U(1)$ subgraph is a \emph{forest} with no loops.
    \item There are \textbf{zero 2-step paths} through $U(1)$ edges. A fermion reached via a $U(1)$ edge has no further $U(1)$ neighbors to propagate through.
    \item Each Type-A vertex (gauge boson) has exactly 6 $U(1)$ edges in its neighborhood, uniformly distributed.
    \item All $U(1)$ pairs are cross-coset: for each pair of cosets, exactly 4 $U(1)$ edges connect fermions from those two cosets.
\end{itemize}

\textbf{Spectral signature.} The $U(1)$ Laplacian (restricted to the 48-edge subgraph on 96 vertices) has a single non-zero eigenvalue: $\lambda = 2$ with multiplicity 48. The spectral gap is $\Delta = 2$---the \emph{maximum possible} for a graph of degree 1.

\textbf{Physical interpretation.} Confinement of the $U(1)$ sector is \emph{inevitable}: with degree 1, no propagation paths exist beyond the immediate partner. This is not a dynamical phenomenon requiring energy scales or running couplings---it is a topological consequence of the edge structure. The photon interacts with charged fermions only through isolated pair correlations, with no extended $U(1)$ flux tubes or Wilson loops.

\textbf{Status: DERIVED.} The perfect matching property, zero cycle rank, and maximum spectral gap follow directly from the edge classification (Section~\ref{sec:edge_split}).


\subsection{Gauge Channel Hierarchy (EXP-164)}
\label{sec:gauge_hierarchy}

The three gauge channels exhibit a characteristic hierarchy in their propagation ranges, determined by the shared-neighbor structure.

\textbf{Propagation range.} Define the range of a gauge channel as the maximum graph distance at which same-channel 2-step paths exist:
\begin{center}
\begin{tabular}{lccc}
\toprule
\textbf{Channel} & \textbf{Edges} & \textbf{Range} & \textbf{Amplitude at $d=1$} \\
\midrule
$U(1)$ & 48 & 1 (ultralocal) & 4.76 \\
$SU(2)_{\text{AC}}$ & 96 & 2 & 2.38 \\
$SU(2)_{\text{BC}} + SU(3)$ & $192 + 384$ & $\geq 5$ (long-range) & 0.79 \\
\bottomrule
\end{tabular}
\end{center}

The hierarchy $\text{range}_{U(1)} < \text{range}_{SU(2)} < \text{range}_{SU(3)}$ matches the physical expectation: the electromagnetic interaction is the most ``local'' (confined to pairs, Section~\ref{sec:u1_confinement}), the weak interaction has intermediate range, and the strong interaction is long-range on the graph (reaching the full diameter).

\textbf{Box diagrams.} The 4-point (box) interaction vertices total 288, decomposing as $192\,(AB) + 96\,(BB) + 0\,(AA)$. The absence of pure gauge-boson box diagrams ($AA = 0$) is a direct consequence of the selection rule A-A $= 0$ (no edges between gauge vertices).

\textbf{Three-point ratio.} The ratio of Yukawa (B-mediated) to gauge (A-mediated) 3-point vertices is $B/A = 16/8 = 2$, reflecting the Type-B to Type-A vertex count ratio.

\textbf{Status: DERIVED.} The propagation ranges, box diagram counts, and vertex ratios follow from the graph structure.


%==============================================================
\section{Coupling Constants and Higgs Mass}
%==============================================================

\subsection{The Fine Structure Constant}

The fine structure constant $\alpha$ emerges from a consistency equation relating the Laplacian spectrum to geometric phases:

\begin{equation}
\boxed{2\pi\alpha^2 - 20\phi^4\alpha + 1 = 0}
\end{equation}

\textbf{Origin of terms:}
\begin{itemize}
    \item $20\phi^4 = 5 \times (\lambda_4/\lambda_1)^2$ from the 600-cell Laplacian spectrum
    \item $\lambda_4/\lambda_1 = 2\phi^2$ (exact eigenvalue ratio)
    \item $2\pi$ = geometric phase accumulated on a Hopf fiber (decagon loop)
\end{itemize}

\textbf{Solution:}
\begin{equation}
\alpha = \frac{20\phi^4 - \sqrt{(20\phi^4)^2 - 8\pi}}{4\pi} = \frac{1}{137.036...}
\end{equation}

\begin{table}[H]
\centering
\begin{tabular}{lcc}
\toprule
& \textbf{Calculated} & \textbf{Experimental \cite{codata2022}} \\
\midrule
$\alpha^{-1}$ & 137.036 & $137.035999084(21)$ \\
\textbf{Error} & \multicolumn{2}{c}{\textbf{0.0001\%}} \\
\bottomrule
\end{tabular}
\caption{Fine structure constant comparison.}
\end{table}

\subsection{The Strong Coupling Constant}\label{sec:alpha_s}

The strong coupling constant at the $Z$ mass scale:

\begin{equation}
\alpha_s(M_Z) = \frac{1}{2\phi^3} = 0.1180
\end{equation}

\textbf{Exact algebraic relation:}
\begin{equation}
\boxed{\frac{\alpha_s}{\alpha} = 10\phi}
\end{equation}

This ratio has a spectral derivation. From Table~\ref{tab:spectrum}, the eigenvalue ratio $\theta_1/\theta_3 = 6\phi/3 = 2\phi$, and the intersection number $a_1 = 5$ (Section~\ref{sec:spectral}), so:
\begin{equation}
\frac{\alpha_s}{\alpha} = a_1 \cdot \frac{\theta_1}{\theta_3} = 5 \times 2\phi = 10\phi
\end{equation}

\textbf{Classification:} DERIVED---the ratio $\alpha_s/\alpha$ is determined by the first intersection number ($a_1 = 5$) and an exact eigenvalue ratio ($\theta_1/\theta_3 = 2\phi$), both topological invariants of the 600-cell.

\begin{table}[H]
\centering
\begin{tabular}{lcc}
\toprule
& \textbf{Calculated} & \textbf{Experimental} \\
\midrule
$\alpha_s(M_Z)$ & 0.1180 & $0.1179 \pm 0.0009$ \\
\textbf{Error} & \multicolumn{2}{c}{\textbf{0.03\%}} \\
\bottomrule
\end{tabular}
\caption{Strong coupling constant comparison.}
\end{table}

\subsection{The Weinberg Angle}

The electroweak mixing angle derives from geometric structure counts:

\begin{equation}
\sin^2\theta_W = \frac{6}{26} = \frac{6}{6+20} = 0.2308
\end{equation}

\textbf{Geometric interpretation:}
\begin{itemize}
    \item 6 = decagonal directions per vertex (U(1) sector)
    \item 20 = tetrahedral cells per vertex (SU(2) sector)
    \item 26 = total structures
\end{itemize}

\textbf{Second derivation path (spectral):} The same ratio emerges from the association scheme (Section~\ref{sec:spectral}):
\begin{equation}
\sin^2\theta_W = \frac{b_1}{\dim(\phi\text{-sector})} = \frac{6}{26}
\end{equation}
where $b_1 = 6$ is the second intersection number and $\dim(\phi\text{-sector}) = 4 + 9 + 9 + 4 = 26$ is the total dimension of the four $\phi$-dependent eigenspaces. Two independent geometric arguments yielding the same fraction strengthens the derivation.

\begin{table}[H]
\centering
\begin{tabular}{lcc}
\toprule
& \textbf{Calculated} & \textbf{Experimental} \\
\midrule
$\sin^2\theta_W$ & 0.2308 & 0.2312 \\
\textbf{Error} & \multicolumn{2}{c}{\textbf{0.19\%}} \\
\bottomrule
\end{tabular}
\caption{Weinberg angle comparison.}
\end{table}


\subsection{Mass from Dihedral Angles}

The Higgs mass represents the phase transition energy from tetrahedral (600-cell) to octahedral (24-cell) geometry:

\textbf{Bare mass:}
\begin{equation}
m_H^{(0)} = m_W \cdot \frac{\theta_{\text{octahedron}}}{\theta_{\text{tetrahedron}}} = m_W \cdot \frac{\arccos(-1/3)}{\arccos(1/3)} = 124.76 \text{ GeV}
\end{equation}

\textbf{Topological correction:}
The 8 vertices of the inscribed 16-cell contribute a correction of $-8\alpha$:

\begin{equation}
\boxed{m_H = m_W \cdot (\phi - 8\alpha) = 125.36 \text{ GeV}}
\end{equation}

\begin{table}[H]
\centering
\begin{tabular}{lcc}
\toprule
& \textbf{Calculated} & \textbf{Experimental (LHC)} \\
\midrule
$m_H$ & 125.36 GeV & 125.25 GeV \\
\textbf{Error} & \multicolumn{2}{c}{\textbf{0.09\%}} \\
\bottomrule
\end{tabular}
\caption{Higgs mass comparison.}
\end{table}


%==============================================================
\section{Fermion Content and Generations}
%==============================================================

\subsection{Three Generations from Geometry}

The 96 vertices of the snub 24-cell (golden-ratio vertices) accommodate exactly three generations of fermions:

\begin{equation}
96 = 16 \times 3 \times 2
\end{equation}

\begin{itemize}
    \item 16 = Weyl fermions per generation (quarks + leptons with color)
    \item 3 = generations
    \item 2 = chiralities (L/R)
\end{itemize}

\subsection{The 4-to-3 Generation Selection Mechanism}

A deeper analysis (EXP-090) reveals that the 96 vertices naturally decompose into \textbf{four groups of 24}, based on which coordinate vanishes:
\begin{equation}
96 = 4 \times 24
\end{equation}

\textbf{Key observation:} Each axial direction $k$ in $\mathbb{R}^4$ connects only to the groups where $x_k \neq 0$. Therefore:

\begin{itemize}
    \item Choosing a \textbf{privileged axis} (the time direction, e.g., $x_3$) \textbf{selects 3 groups}
    \item The fourth group (where $x_3 = 0$) becomes \textbf{disconnected} from the temporal axis
    \item This group is identified with a \textbf{sterile sector} (right-handed neutrinos, dark matter candidates)
\end{itemize}

\textbf{Physical interpretation:} The breaking of the $\mathbb{R}^4$ symmetry by the choice of time direction naturally reduces 4 geometric groups to 3 observable generations plus 1 sterile sector.

\subsection{Spectral Derivation of Three Generations (EXP-180, 181)}
\label{sec:three_splitting}

A second, independent derivation of the number three arises from the spectral perturbation theory of topological solitons on the 600-cell graph.

\textbf{Setup.} Consider a single soliton defect: one vertex $v_0$ is flipped from its ground-state coset to a different coset (Section~\ref{sec:soliton_creation}). This modifies the graph Laplacian $L \to L'$ by removing the 3 edges connecting $v_0$ to neighbors in the target coset and adding 3 new edges. The perturbation has rank 3.

\textbf{Universal splitting pattern.} Each of the 9 eigenvalue groups of the unperturbed 600-cell Laplacian (with multiplicities $m_k \in \{1, 4, 9, 16, 25, 36, 9, 16, 4\}$) splits into exactly \textbf{three sub-levels} under the soliton perturbation, with degeneracies:
\begin{equation}
\boxed{m_k \;\to\; (m_k - 3,\; 2,\; 1)}
\end{equation}

This pattern is \textbf{universal}: it is independent of the choice of vertex $v_0$, its type (A, B, or C), and the target coset. The number 3 has a geometric origin:
\begin{equation}
3 = \frac{12}{4} = \frac{\text{neighbors per vertex}}{\text{other cosets}}
\end{equation}
Each vertex has 12 neighbors distributed equally among 4 other cosets (3 per coset), and the soliton perturbs exactly the 3 edges to the target coset.

\textbf{Exact eigenvector structure (EXP-181).} Within each eigenspace $\lambda_k$ of multiplicity $m_k$, the three split sub-levels have distinct spatial character. Denoting the projector $P_k$ restricted to the soliton support (vertex $v_0$ and its 12 neighbors), the perturbation matrix $M = P_k \Delta L P_k$ has eigenvalues $\{0, 1, 1, 4\}$ (on the 4-dimensional support subspace), giving shifts:

\begin{itemize}
    \item \textbf{Bulk mode} ($m_k - 3$ degenerate): unshifted, $\Delta\lambda = 0$. Eigenvectors vanish on the soliton support.
    \item \textbf{Shell doublet} (2-fold): shift $\Delta\lambda = \lambda_k - \tau_k$, where $\tau_k$ is the off-diagonal projector element. The doublet eigenvectors satisfy $|\psi(v_0)|^2 = 0$ \emph{exactly}---they live \textbf{only} on the neighbor shell.
    \item \textbf{Core singlet} (1-fold): shift $\Delta\lambda = 4\sigma_k$, where $\sigma_k$ is the diagonal projector element. The singlet is localized at the defect, with $|\psi(v_0)|^2$ increasing with eigenvalue.
\end{itemize}

\textbf{Physical interpretation.} The three sub-levels correspond to three distinct localization patterns around a topological defect:
\begin{center}
\begin{tabular}{lcl}
\textbf{Sub-level} & \textbf{Localization} & \textbf{Generation analogy} \\
\hline
Bulk ($m_k - 3$) & Delocalized, avoids defect & 1st generation (light: electron) \\
Shell doublet (2) & On neighbor shell only & 2nd generation (intermediate: muon) \\
Core singlet (1) & At defect site & 3rd generation (heavy: tau) \\
\end{tabular}
\end{center}

This matches the Yukawa hierarchy: the core-localized mode couples most strongly to the soliton (heaviest), while the bulk mode is insensitive (lightest). The doublet's exact vanishing at $v_0$ provides a selection rule separating the second generation from the third.

\textbf{Uniqueness.} A control analysis shows that the uniform 3-way splitting occurs \textbf{only} when exactly 3 edges are perturbed (= one soliton). Perturbing 1 or 2 edges gives only 2 sub-levels; perturbing 4 or more edges breaks the universality (results depend on vertex choice). Thus the 600-cell graph structure \emph{uniquely selects} a 3-fold splitting.

\textbf{Status:} The number of sub-levels (three) is \textbf{DERIVED} from the perturbation rank and the uniform coset structure. The identification of sub-levels with fermion generations is a \textbf{PATTERN} (the localization hierarchy matches the mass hierarchy qualitatively, but the shift ratios $v_3/v_2$ are not constant across eigenspaces).


\subsection{The Generation Theorem: $N_{\text{gen}} = 3$ (EXP-172, 182--186)}
\label{sec:fibonacci_generations}

We prove that the number of fermion generations is exactly three. The proof combines the algebraic structure of $\mathbb{Z}[\phi]$ with the Galois invariance of $E_8$.

\begin{quote}
\textbf{Theorem.} \textit{On the $a=1$ mass line, exactly three values of $b$ yield stable fermion states: $b \in \{0, 1, 2\}$. The number three is an immutable property of the Fibonacci sequence and cannot be adjusted by any continuous parameter.}
\end{quote}

\textbf{Step 1 (Geometry $\to$ Algebra).} The 600-cell adjacency spectrum lies in $\mathbb{Z}[\phi] = \mathbb{Z} + \mathbb{Z}\phi$. The intersection numbers $a_1 = 5$, $b_1 = 6$ define the mass level $n = 5a + 6b$ and the algebraic element $z = a + b\phi$ (Section~\ref{sec:dsi}).

\textbf{Step 2 ($E_8$ $\to$ Galois).} The $E_8$ icosian construction (Section~\ref{sec:e8}) decomposes $E_8 = S \cup T$ where $T = \phi' \cdot S$. A fermion at $z = a + b\phi$ in the $S$-sector has a \emph{Galois shadow} $z' = \sigma(z) = a + b\phi'$ in the $T$-sector, with algebraic norm $N(z) = z \cdot z' = a^2 + ab - b^2 \in \mathbb{Z}$.

\textbf{Step 3 (Galois Invariance $\to$ Unit Selection).} The DSI potential $V(\psi) = \sin^2(\pi \ln|\psi| / \ln\phi)$ satisfies
\begin{equation}
\sigma(V(z)) = V(z')
\label{eq:galois_covariance}
\end{equation}
since $\sigma(\phi) = \phi'$, $\ln|\phi'| = -\ln\phi$, and $\sin^2(-x) = \sin^2(x)$. The $E_8$ lattice is defined over $\mathbb{Z}$: its Gram matrix has integer entries, and physical observables derived from $E_8$ must be Galois-invariant (independent of the choice $\sqrt{5} \leftrightarrow -\sqrt{5}$). Since $V(z)$ alone is \emph{not} Galois-invariant for $b \geq 4$, the minimal Galois-invariant energy functional is
\begin{equation}
E(z) = V(z) + V(z') = \operatorname{Tr}_{\mathbb{Q}(\phi)/\mathbb{Q}}\bigl(V(z)\bigr)
\label{eq:galois_energy}
\end{equation}
which is the field trace---the canonical Galois-invariant extension. Since $V \geq 0$, stability ($E = 0$) requires $V(z) = 0$ \emph{and} $V(z') = 0$, i.e., $|z| = \phi^m$ and $|z'| = \phi^n$ for integers $m, n$. Then $|N(z)| = \phi^{m+n}$. But $N(z) \in \mathbb{Z}$, and $\phi^k \in \mathbb{Q}$ only for $k = 0$. Therefore $|N(z)| = 1$: the element $z$ is a \emph{unit} of $\mathbb{Z}[\phi]$.

\textbf{Step 4 (Dirichlet $\to$ Fibonacci).} By Dirichlet's unit theorem, the units of $\mathbb{Z}[\phi]$ are $\{\pm\phi^k : k \in \mathbb{Z}\}$. Since $\phi^k = F(k\!-\!1) + F(k)\phi$ where $F$ denotes the Fibonacci sequence, the condition $z = 1 + b\phi$ (i.e., the $a = 1$ line) requires $F(k\!-\!1) = 1$. The equation $F(m) = 1$ has exactly three solutions: $m \in \{-1, 1, 2\}$ (for $m \geq 3$, $F(m) \geq 2$ and the Fibonacci sequence is strictly increasing). This gives $k \in \{0, 2, 3\}$ and therefore $b = F(k) \in \{0, 1, 2\}$.

\textbf{Step 5 (Conclusion).} Exactly three generations on the $a = 1$ line:
\begin{center}
\begin{tabular}{ccccl}
\toprule
$b$ & $k$ & $z = \phi^k$ & $n = 5+6b$ & Fermion \\
\midrule
0 & 0 & $\phi^0 = 1$ & 5 & electron / down \\
1 & 2 & $\phi^2 = 1+\phi$ & 11 & muon / strange \\
2 & 3 & $\phi^3 = 1+2\phi$ & 17 & tau / bottom \\
\bottomrule
\end{tabular}
\end{center}
For $b = 3$: $z = 1 + 3\phi \neq \phi^k$ for any integer $k$, and $|N| = 5 \neq 1$, so $E(z) > 0$. \hfill $\square$

\textbf{Triple identity.} The number $a_1 = 5$ plays three simultaneous roles: it is the graph diameter, the discriminant $\operatorname{disc}(\mathbb{Z}[\phi]) = 5$, and the number of cosets $|2I/2T| = 5$. This unifies the graph-theoretic, algebraic, and group-theoretic routes to $N_{\text{gen}} = 3$.

\textbf{Non-unit fermions.} For $b \geq 3$ on $a = 1$, $|N|$ grows: $|N(1,3)| = 5$, $|N(1,4)| = 11$. The heavy quarks charm ($|N| = 5$), bottom ($|N| = 19$), and top ($|N| = 19$) are non-units. The top quark, with the largest $|N|$, is the only quark decaying before hadronizing---a suggestive correlation. Of the 9 fermions, 6 are units and 3 are non-units; the non-units are always the heaviest in their family.

\textbf{Status:} All five steps use established mathematics (Galois theory, Dirichlet's unit theorem, Fibonacci properties). The axioms are: (A1) the 600-cell as starting geometry, (A2) $E_8$ as ambient lattice, (A3) DSI potential from holonomy. From these, $N_{\text{gen}} = 3$ follows without free parameters. \textbf{Classification: DERIVED.}


\subsection{Representation Branching and Bare Flavor Mixing (EXP-215--217b)}
\label{sec:branching_mixing}

The binary icosahedral group $2I$ has 9 irreducible representations with dimensions $\{1, 2, 2, 3, 3, 4, 4, 5, 6\}$, all with characters in $\mathbb{Q}(\phi)$. The binary tetrahedral subgroup $2T \subset 2I$ has 7 irreducible representations with dimensions $\{1, 1, 1, 2, 2, 2, 3\}$. Crucially, $2T$ has \textbf{three} 2-dimensional irreps: $\psi_5$ (real) and $\psi_3, \psi_4$ (complex conjugate pair with $\omega = e^{2\pi i/3}$ characters).

\textbf{Key branching rule.} The 6-dimensional irrep $\rho_8$ of $2I$ (corresponding to the eigenvalue $\lambda_5 = -2$ with multiplicity $36 = 6^2$) branches as:
\begin{equation}
\boxed{\rho_8|_{2T} = \psi_3 \oplus \psi_4 \oplus \psi_5 \quad \text{(all three 2D irreps of } 2T\text{)}}
\end{equation}
In contrast, the branch node representation (dim~4, $\lambda_3 = 3$) branches as $\psi_3 \oplus \psi_4$ (complex pair only, no real component), and the $E_8$ leg endpoint representations (dim~2) restrict to $\psi_5$ (real 2D only). This provides a \textbf{structural asymmetry} between the three generations at the representation level.

\textbf{Conjugate subgroup structure.} Since $A_5$ is simple, $2T$ is \emph{not} a normal subgroup of $2I$, and there are 5 distinct conjugate subgroups $2T_g = g \cdot 2T \cdot g^{-1}$---one for each coset. This matches $|2I/2T| = 5 = a_1$.

\textbf{Democratic bare mixing (EXP-217b).} Using the unitary Schur projectors for the three 2D irreps, the mixing matrix between any pair of conjugate subgroups $(2T_g, 2T_h)$ for $g \neq h$ is:
\begin{equation}
\boxed{M = \begin{pmatrix} 1/6 & 5/12 & 5/12 \\ 5/12 & 1/6 & 5/12 \\ 5/12 & 5/12 & 1/6 \end{pmatrix}}
\label{eq:bare_mixing}
\end{equation}
This matrix is doubly stochastic and $S_3$-symmetric. The diagonal element $p = 1/6$ is \textbf{derived} from Schur's lemma:
\begin{equation}
p = \frac{n \cdot r - 1}{n - 1}, \qquad n = 5 \;(= |2I/2T|), \quad r = \frac{d_{\text{irrep}}}{d_{\text{rep}}} = \frac{2}{6} = \frac{1}{3}
\end{equation}
The 2-transitivity of $2I$ on $2I/2T$ ensures that \emph{all 10 pairwise mixing matrices are identical}---a consequence of the simplicity of $A_5$.

\textbf{Physical interpretation.} The bare mixing is fully democratic: no preferred direction in generation space. Physical CKM and PMNS mixing matrices require \textbf{mass-induced symmetry breaking} to lift this $S_3$ degeneracy. The $\rho_8$ branching into $\psi_3 \oplus \psi_4 \oplus \psi_5$ provides the necessary structural asymmetry (two complex + one real), while the mass formula (Section~6.6) provides the dynamical breaking.

\textbf{Status:} The branching rules and $p = 1/6$ are \textbf{DERIVED}. The democratic bare mixing is \textbf{DERIVED}. The connection to physical CKM/PMNS requires mass-breaking mechanisms (OPEN).

%==============================================================
\section{Fermion Mass Hierarchy from Holonomy}
%==============================================================

\subsection{The Mass Ratio Pattern}

The lepton mass ratios follow a remarkable pattern based on powers of $\phi$:

\begin{table}[H]
\centering
\begin{tabular}{lcccc}
\toprule
\textbf{Ratio} & \textbf{$n$} & \textbf{$\phi^n$} & \textbf{Experimental} & \textbf{Error} \\
\midrule
$m_\mu/m_e$ & 11 & 199.0 & 206.77 & 3.8\% \\
$m_\tau/m_\mu$ & 6 & 17.94 & 16.82 & 6.7\% \\
$m_\tau/m_e$ & 17 & 3571.0 & 3477.1 & 2.7\% \\
\bottomrule
\end{tabular}
\caption{Lepton mass ratios as powers of $\phi$.}
\end{table}

\subsection{Derivation from 600-Cell Geometry}

The exponents 5, 6, 11, 17 are \textbf{derived} from the intrinsic geometry of the 600-cell:

\begin{equation}
\boxed{
\begin{aligned}
5 &= \text{diameter of the 600-cell graph} \\
6 &= \text{decagons per vertex} \\
11 &= 5 + 6 \quad \text{(holonomy additivity)} \\
17 &= 11 + 6 \quad \text{(second excitation level)}
\end{aligned}
}
\end{equation}

\subsection{The Hopf Fibration Connection}

The 600-cell, inscribed in $S^3$, inherits the \textbf{Hopf fibration} structure $S^3 \to S^2$:

\begin{itemize}
    \item 120 vertices $\to$ 30 fibers $\times$ 4 vertices each
    \item Edge angle $= \pi/5 = 36°$ \textbf{(exact)}
    \item Total phase on decagon $= 2\pi$
\end{itemize}

\textbf{The holonomy argument:} When a spinor field is parallel-transported along a path on $S^3$:
\begin{itemize}
    \item Traversing the \textbf{diameter} (5 steps) accumulates phase $\sim 5 \cdot \log\phi$
    \item Each \textbf{decagon} (6 per vertex) contributes phase $\sim 1 \cdot \log\phi$
    \item Total for muon: $(5 + 6) \cdot \log\phi = 11 \cdot \log\phi$
\end{itemize}

The mass formula becomes:
\begin{equation}
m = m_e \cdot \phi^n = m_e \cdot e^{n \log\phi}
\end{equation}

\subsection{The Binary Icosahedral Group}

The holonomy group of the 600-cell is related to the \textbf{binary icosahedral group} $2I$:
\begin{itemize}
    \item $|2I| = 120$ (same as number of vertices!)
    \item Irreducible representations involve powers of $\phi$
    \item Holonomy is ``quantized'' in units related to $\phi$
\end{itemize}

\textbf{Conclusion:} The exponents 5, 6, 11, 17 are not numerological coincidences---they are \textbf{topological invariants} of the 600-cell lattice structure.

\subsection{Electroweak Phase Corrections (EXP-097)}

The bare exponents $n=11,\,6,\,17$ yield errors of 3.8\%, 6.7\%, and 2.7\% respectively. We decompose the effective exponents into their geometric components:

\begin{equation}
n_\mu = d_{\text{eff}} + k_{\text{eff}}, \quad
n_{\tau/\mu} = k_{\text{eff}}, \quad
n_\tau = d_{\text{eff}} + 2\,k_{\text{eff}}
\end{equation}

where $d_{\text{eff}} = 5 + \delta_d$ (diameter contribution) and $k_{\text{eff}} = 6 + \delta_k$ (decagonal contribution). A systematic numerical search (EXP-097) identifies:

\begin{equation}
\boxed{
\delta_d = \sin^2\theta_W = \frac{6}{26}, \qquad
\delta_k = -\frac{1}{\phi^4}
}
\end{equation}

\textbf{Physical interpretation:}
\begin{itemize}
    \item The diameter path acquires a correction proportional to the electroweak mixing angle --- the same ratio $6/26$ already derived from vertex structure counts.
    \item The decagonal loop acquires a suppression of $1/\phi^4$, reflecting the four-dimensional embedding of the geodesic.
\end{itemize}

The corrected exponents give:

\begin{table}[H]
\centering
\begin{tabular}{lcccc}
\toprule
\textbf{Ratio} & \textbf{$n_{\text{eff}}$} & \textbf{Calculated} & \textbf{Experimental} & \textbf{Error} \\
\midrule
$m_\mu/m_e$ & 11.085 & 205.95 & 206.77 & \textbf{0.28\%} \\
$m_\tau/m_e$ & 16.939 & 1772.4 & 1776.9 & \textbf{0.25\%} \\
$m_\tau/m_\mu$ & 5.854 & 16.73 & 16.82 & \textbf{0.53\%} \\
\bottomrule
\end{tabular}
\caption{Lepton mass ratios with electroweak phase corrections.}
\end{table}

\textbf{Status:} The corrections $\sin^2\theta_W$ and $1/\phi^4$ are quantities already present in the theory (Sections 3.3 and 2.1), not free parameters. This reduces the lepton mass errors from 3--7\% to below 0.5\%.

\subsection{Unified Fermion Mass Formula (EXP-098--100)}

The electroweak phase corrections generalize to all fermions through sector-dependent corrections. The unified mass formula reads:

\begin{equation}
\boxed{m_f = m_e \cdot \phi^{\,a \cdot d_{\text{eff}} \;+\; b \cdot k_{\text{eff}}}}
\end{equation}

where $(a,b) \in \mathbb{Z}^2$ are \textbf{topological quantum numbers} specific to each fermion, and:
\begin{equation}
d_{\text{eff}} = 5 + \delta_d, \qquad k_{\text{eff}} = 6 + \delta_k
\end{equation}

The corrections $\delta_d, \delta_k$ depend on the fermion sector and reflect the dominant interaction:

\begin{table}[H]
\centering
\begin{tabular}{lccc}
\toprule
\textbf{Sector} & $\delta_d$ & $\delta_k$ & \textbf{Interpretation} \\
\midrule
Leptons ($Q=-1$) & $+\sin^2\theta_W$ & $-1/\phi^4$ & Electroweak \\
Up-type ($Q=+\tfrac{2}{3}$) & $+2\alpha_s/\pi$ & $+\alpha_s$ & QCD (positive shift) \\
Down-type ($Q=-\tfrac{1}{3}$) & $-1/5$ & $-\alpha_s$ & QCD (negative shift) \\
\bottomrule
\end{tabular}
\caption{Sector-dependent corrections. All quantities derive from the theory (Sections~4.2, 4.3).}
\end{table}

\textbf{Key point:} No new parameters are introduced. The lepton corrections use $\sin^2\theta_W$ (Section~4.3) and $\phi^{-4}$ (geometric). The quark corrections use $\alpha_s = 1/(2\phi^3)$ (Section~4.2) and $1/5 = 1/\text{diameter}$ (Section~2.1). The down-type correction $\delta_k = -\alpha_s = -1/(2\phi^3)$ is the same relation already established in Eq.~(6).

The complete fermion mass spectrum:

\begin{table}[H]
\centering
\begin{tabular}{lcccccr}
\toprule
\textbf{Fermion} & $Q$ & $(a,b)$ & $n$ & $m_{\text{calc}}$ & $m_{\text{exp}}$ & \textbf{Error} \\
\midrule
Electron & $-1$ & $(0,0)$ & 0 & 0.511 MeV & 0.511 MeV & 0.00\% \\
Muon & $-1$ & $(1,1)$ & 11 & 106.0 MeV & 105.7 MeV & \textbf{0.28\%} \\
Tau & $-1$ & $(1,2)$ & 17 & 1772 MeV & 1777 MeV & \textbf{0.25\%} \\
\midrule
Up & $+\tfrac{2}{3}$ & $(3,-2)$ & 3 & 2.15 MeV & 2.16 MeV & \textbf{0.30\%} \\
Charm & $+\tfrac{2}{3}$ & $(2,1)$ & 16 & 1283 MeV & 1270 MeV & \textbf{1.03\%} \\
Top & $+\tfrac{2}{3}$ & $(4,1)$ & 26 & 169.6 GeV & 172.8 GeV & 1.82\% \\
\midrule
Down & $-\tfrac{1}{3}$ & $(1,0)$ & 5 & 5.15 MeV & 4.67 MeV & $10.2\%^\dagger\!^\ddagger$ \\
Strange & $-\tfrac{1}{3}$ & $(1,1)$ & 11 & 87.3 MeV & 93.4 MeV & 6.6\% \\
Bottom & $-\tfrac{1}{3}$ & $(-1,4)$ & 19 & 4191 MeV & 4180 MeV & \textbf{0.26\%} \\
\bottomrule
\end{tabular}
\caption{Complete fermion mass predictions. $n = 5a + 6b$ is the bare exponent. $^\dagger$Down quark experimental mass: $4.67^{+0.48}_{-0.17}$ MeV (PDG~\cite{pdg2024}). $^\ddagger$The prediction 5.15\,MeV coincides with the 1$\sigma$ upper bound ($4.67 + 0.48 = 5.15$\,MeV), placing it within $1\sigma$ of the central value (EXP-128).}
\end{table}

Seven of nine fermions are predicted within 2\% of experiment. The two outliers---Down ($10\%$) and Strange ($7\%$)---are the lightest quarks, where non-perturbative QCD effects and experimental mass uncertainties are largest. A comprehensive investigation (EXP-128) shows that the Down quark prediction (5.15\,MeV) coincides exactly with the PDG 1$\sigma$ upper bound, placing it within $1\sigma$ of the experimental value. The tension between Down and Strange is structural: both lie on the $a = 1$ line in the $(a,b)$ lattice and share the same $\delta_d$ correction, so improving one necessarily worsens the other.

\textbf{Structural patterns:}
\begin{enumerate}
\item \textbf{Strange--Muon duality:} Both share $(a,b) = (1,1)$ and $n=11$, differing only in their sector corrections. At the QCD scale $\mu \approx 1.17$ GeV, their running masses coincide.
\item \textbf{Topological complexity:} The average $|a|+|b|$ increases with generation: Gen~1 (2.0) $<$ Gen~2 (2.3) $<$ Gen~3 (4.3), suggesting heavier fermions require more complex paths on the 600-cell graph.
\item \textbf{Self-consistency:} All correction terms ($\sin^2\theta_W$, $\alpha_s$, $1/\phi^4$, $1/5$) were independently derived elsewhere in the theory.
\end{enumerate}

\subsection{Derivation of $(a,b)$ from Complexity Minimization}
\label{sec:ab_derivation}

The topological quantum numbers $(a,b)$ are \emph{not} free parameters. They are uniquely determined by the 600-cell geometry plus a single principle.

\textbf{Theorem.} For each fermion with mass level $n$, the quantum numbers $(a,b)$ are the \emph{unique} solution minimizing topological complexity:
\begin{equation}
\boxed{(a,b) = \arg\min_{a_1 a + b_1 b = n} \bigl(|a| + |b|\bigr)}
\end{equation}
where $a_1 = 5$ and $b_1 = 6$ are the intersection numbers of the 600-cell (Section~\ref{sec:spectral}).

\textbf{Proof sketch.} The general integer solution of $5a + 6b = n$ is $a = a_0 + 6t$, $b = b_0 - 5t$ for $t \in \mathbb{Z}$. The complexity $f(t) = |a_0 + 6t| + |b_0 - 5t|$ is a convex piecewise-linear function with kinks at $t = -a_0/6$ and $t = b_0/5$. For $n \neq 0$ these kinks are at distinct positions, so $f(t)$ has a unique integer minimizer. $\square$

This has been verified computationally for all 9 fermions:

\begin{table}[H]
\centering
\begin{tabular}{lcccccc}
\toprule
\textbf{Fermion} & $n$ & \textbf{Predicted} $(a,b)$ & $|a|+|b|$ & \textbf{Unique?} & \textbf{Gap} \\
\midrule
Electron & 0 & $(0,0)$ & 0 & Yes & 11 \\
Down & 5 & $(1,0)$ & 1 & Yes & 9 \\
Muon/Strange & 11 & $(1,1)$ & 2 & Yes & 9 \\
Charm & 16 & $(2,1)$ & 3 & Yes & 7 \\
Tau & 17 & $(1,2)$ & 3 & Yes & 7 \\
Up & 3 & $(3,-2)$ & 5 & Yes & 1 \\
Bottom & 19 & $(-1,4)$ & 5 & Yes & 1 \\
Top & 26 & $(4,1)$ & 5 & Yes & 3 \\
\bottomrule
\end{tabular}
\caption{Complexity minimization predicts all 9 fermion $(a,b)$ uniquely. The ``Gap'' column shows $|a|+|b|$ of the next-best decomposition minus the minimum.}
\label{tab:complexity}
\end{table}

The result is robust: $L^1$, $L^2$, and $L^\infty$ norms all select the same $(a,b)$ for all 9 fermions. This reduces the free parameter count from 18 (nine fermions $\times$ two quantum numbers) to 9 (the mass levels $n$ alone).

\textbf{Status:} The quantum numbers $(a,b)$ are DERIVED from complexity minimization with intersection numbers $a_1 = 5$, $b_1 = 6$. The sector-dependent corrections are DERIVED from quantities already in the theory. The remaining input is the set of 9 mass levels $n$, which come from the DSI mechanism (Section~\ref{sec:dsi}) and experimental masses. The Down quark ($10\%$) remains an open problem.

\subsection{Galois Complementarity of Mass Corrections (EXP-187--189)}
\label{sec:galois_complementarity}

The sector-dependent corrections $\delta_d$, $\delta_k$ to the effective exponents $d_{\text{eff}} = 5 + \delta_d$ and $k_{\text{eff}} = 6 + \delta_k$ reveal a striking algebraic pattern when expressed in $\mathbb{Z}[\phi]$ coordinates. Writing $\delta z = \delta_d + \delta_k \cdot \phi$ and $\delta z' = \delta_d + \delta_k \cdot \phi'$:

\begin{table}[H]
\centering
\begin{tabular}{lccccl}
\toprule
\textbf{Sector} & $\delta_d$ & $\delta_k$ & $|\delta z|$ & $|\delta z'|$ & \textbf{Pattern} \\
\midrule
Leptons & $+0.214$ & $-0.135$ & $0.004$ & $0.298$ & Shadow only \\
Up quarks & $+0.104$ & $+0.119$ & $0.296$ & $0.031$ & Real only \\
Down quarks & $+0.389$ & $+0.225$ & $0.753$ & $0.249$ & Both sectors \\
\bottomrule
\end{tabular}
\caption{Galois decomposition of mass corrections. For leptons, $|\delta z| / |\delta z'| = 1:78$; for up quarks, $|\delta z'| / |\delta z| = 1:10$. The correction direction is complementary.}
\end{table}

For \textbf{leptons}, the correction nearly vanishes in the $\phi$-sector ($|\delta z| \approx 0$) and lives entirely in the Galois shadow. The correction vector $(\delta_d, \delta_k)$ is proportional to $(\phi, -1)$---the Galois conjugate direction---with $\delta_d / \delta_k = -1.59 \approx -\phi$ (within 2\%).

For \textbf{up quarks}, the pattern is inverted: $\delta_k \approx \alpha_s = 0.1179$ to 0.5\% precision, and $\delta_d \approx \alpha_s(1 - \alpha_s)$ to 0.27\%. The correction lives in the $\phi$-sector with the shadow nearly zero.

This \textbf{Galois complementarity}---leptons corrected in $\phi'$ only, up quarks in $\phi$ only, down quarks in both---is suggestive of a connection between the Galois structure of $\mathbb{Q}(\sqrt{5})$ and the gauge quantum numbers (color charge, electric charge). However, a derivation from first principles has not yet been found.

\textbf{Near-exact Galois flip (EXP-190).} The large components satisfy a striking near-equality:
\begin{equation}
\frac{\delta z'(\text{leptons})}{\delta z(\text{up quarks})} = \frac{0.2977}{0.2960} = 1.0055
\end{equation}
so the lepton shadow correction and the up-quark real correction are equal to $0.55\%$---a near-exact Galois flip between sectors. Furthermore, the small residual of the up-quark shadow satisfies $|\delta z'(\text{up})| \approx \alpha \cdot \phi^3$ to $0.4\%$.

\textbf{Status:} The Galois complementarity is an empirical PATTERN (not derived). The up-quark $\alpha_s$ matching and the near-exact Galois flip are notable numerical coincidences.


%==============================================================
\section{Discrete Scale Invariance: The Mass Mechanism}
\label{sec:dsi}
%==============================================================

The unified mass formula of Section~6.6 raises a natural question: \emph{why} do fermion masses scale as powers of $\phi$? In this section we show that the answer is \textbf{Discrete Scale Invariance} (DSI) \cite{sornette1998}---the mass spectrum arises from a potential with periodicity in logarithmic energy space, and the 600-cell lattice provides both the scaling factor and the selection rules.

\subsection{The DSI Potential}

A scalar field on the 600-cell lattice experiences a self-similar potential with minima at discrete mass scales. The natural form is:

\begin{equation}
\boxed{V(\psi) = A \sin^2\!\left(\frac{\pi \ln(|\psi|/m_e)}{\ln\phi}\right)}
\label{eq:dsi_potential}
\end{equation}

where $A$ sets the barrier height. This potential has degenerate minima at $|\psi| = m_e \cdot \phi^n$ for all integers $n$, defining a \textbf{DSI ladder} of allowed mass levels.

\textbf{Why $\phi$?} The scaling factor is not a free parameter---it is fixed by the 600-cell geometry. The edge angle of the 600-cell inscribed in $S^3$ is $\pi/5$ (exact), and $2\cos(\pi/5) = \phi$. Under parallel transport, each edge contributes a holonomy phase $\ln\phi$ to a spinor field, so a path of $n$ edges yields a mass ratio $e^{n\ln\phi} = \phi^n$.

\subsection{Holonomy as Discrete Scale Invariance}

The equivalence between holonomy on $S^3$ and DSI is exact:

\begin{equation}
\text{Phase per edge} = \ln\phi \quad \Longrightarrow \quad \text{Mass ratio after } n \text{ edges} = \phi^n
\end{equation}

The Hopf fibration $S^3 \to S^2$ decomposes paths on the 600-cell into two fundamental classes:
\begin{itemize}
    \item \textbf{Diameter traversals:} 5 edges across the graph (graph diameter = 5)
    \item \textbf{Decagonal loops:} 6 edges around a decagonal geodesic (6 decagons per vertex)
\end{itemize}

A general path wrapping $a$ times across the diameter and $b$ times around decagonal loops traverses $n = 5a + 6b$ edges, giving:
\begin{equation}
m = m_e \cdot \phi^{5a + 6b}
\end{equation}

The coefficients 5 and 6 have a graph-theoretic origin: they are the \textbf{intersection numbers} $a_1 = 5$ and $b_1 = 6$ of the 600-cell association scheme (Section~\ref{sec:spectral}). Thus $n = a_1 a + b_1 b$, and the mass formula is fully determined by topological invariants.

With sector-dependent corrections (Section~6.5--6.6), this becomes the unified mass formula $m_f = m_e \cdot \phi^{a \cdot d_{\text{eff}} + b \cdot k_{\text{eff}}}$.

\textbf{Physical interpretation:} Fermions are \textbf{topological solitons} sitting at minima of the DSI potential (see Sections~\ref{sec:soliton_creation}--\ref{sec:condensation} for the creation mechanism and energy budget). Each soliton is characterized by its winding numbers $(a,b)$ on the coarse 600-cell graph. The mass is not fitted---it is determined by which potential minimum the soliton occupies. The specific $(a,b)$ for each fermion follows from complexity minimization (Section~\ref{sec:ab_derivation}).

\subsection{Level Selection and Empty Levels}
\label{sec:level_selection}

The DSI ladder has minima at $m_e \cdot \phi^n$ for \emph{all} integers $n = 0, 1, 2, \ldots$. Since $\gcd(5,6)=1$, every non-negative integer $n$ can be written as $5a + 6b$ for some $(a,b) \in \mathbb{Z}^2$. Of the 27 levels from $n=0$ to $n=26$, only 8 are occupied by known fermions (with muon and strange sharing $n=11$).

Two algebraic constraints substantially narrow the allowed set:

\textbf{Condition 1: $\mathbb{Z}[\phi]$ norm constraint.} For each level $n$ with minimal-complexity representative $(a,b)$, compute the algebraic norm $N(a + b\phi) = a^2 + ab - b^2$. Occupied levels have $|N| \in \{0, 1, 5, 19\}$ only---excluding all levels with $|N| \in \{4, 9, 11, 16, 25, \ldots\}$. This reduces 27 candidates to 14.

\textbf{Condition 2: Three-line rule in the $(a,b)$ lattice.} The remaining selection is geometric: an allowed level is occupied only if its $(a,b)$ representative lies on
\begin{equation}
\{\text{origin}\} \;\cup\; \{a = 1\} \;\cup\; \{b = 1\} \;\cup\; \{3a + 2b = 5\}
\end{equation}

The three lines have a remarkable structure:
\begin{itemize}
    \item Line $a = 1$: $d(1,0)$, $\mu/s(1,1)$, $\tau(1,2)$ --- leptons + down quark
    \item Line $b = 1$: $\mu/s(1,1)$, $c(2,1)$, $t(4,1)$ --- heavy quarks
    \item Line $3a+2b = 5$: $u(3,-2)$, $\mu/s(1,1)$, $b(-1,4)$ --- cross-generation
\end{itemize}

All three lines intersect at $(1,1)$ = muon/strange, the triple intersection point. This explains the $\mu$--$s$ degeneracy ($n=11$ shared) as a topological consequence: $(1,1)$ is the unique junction of all fermion lines.

Together, the two conditions correctly classify \textbf{28 of 31 levels} ($n = 0$ to $30$). Three false positives remain: $n = 1\,(-1,1)$, $n = 6\,(0,1)$, $n = 23\,(1,3)$---all on lines $a=1$ or $b=1$ but unoccupied.

\textbf{Per-line refinement (EXP-123).} The three false positives are eliminated by per-line constraints, one of which is derivable:
\begin{itemize}
    \item \textbf{Line $b = 1$:} Require the \emph{signed} norm $N(a + b\phi) = a^2 + ab - b^2 > 0$. On this line, $N = a^2 + a - 1$, which is negative for $a \leq 0$ and positive for $a \geq 1$. This excludes both false positives $n = 1\,(a=-1)$ and $n = 6\,(a=0)$. The occupied values $a = \{1, 2, 4\} = \{2^0, 2^1, 2^2\}$ give norms $|N| = \{1, 5, 19\}$---exactly the three non-zero allowed magnitudes. \textbf{Status: DERIVED} (norm positivity in $\mathbb{Z}[\phi]$).
    \item \textbf{Line $a = 1$:} Require $b \leq 2$, excluding $n = 23\,(b=3)$. This is equivalent to limiting the number of generations to three. \textbf{Status: PATTERN} (physically motivated but not geometrically derived).
    \item \textbf{Line $3a + 2b = 5$:} No additional constraint needed; all three solutions are occupied ($3/3$). \textbf{Status: DERIVED} (self-consistent).
\end{itemize}

With these refinements, the selection rule achieves \textbf{31/31} accuracy. The three-generation limit ($b \leq 2$) remains the sole underived component.

\textbf{Classification:} The $\mathbb{Z}[\phi]$ norm constraint is DERIVED. The three-line rule is a PATTERN. The $b = 1$ line refinement via norm positivity is DERIVED. The $a = 1$ line refinement ($b \leq 2 =$ three generations) is PATTERN---deriving the number of fermion generations from geometry remains one of the deepest open problems in theoretical physics (EXP-127).

\subsection{Fractal Self-Similarity}

The symmetry group of the 600-cell satisfies a suggestive identity:
\begin{equation}
|H_4| = 14{,}400 = 120^2
\end{equation}

If interpreted as a \textbf{fractal hierarchy}---each of the 120 vertices of a coarse 600-cell resolving into a finer 600-cell of 120 vertices---then:
\begin{itemize}
    \item \textbf{Level 0} (coarse graph): determines the coupling constants ($\alpha$, $\alpha_s$, $\sin^2\theta_W$) from spectral properties
    \item \textbf{Level 1} (fine structure): determines fermion masses from soliton paths $(a,b)$ on the coarse graph
    \item Sector corrections ($\delta_d$, $\delta_k$) encode the coupling between levels
\end{itemize}

The number of DSI levels from the Planck scale to the electron mass is $\ln(m_P/m_e)/\ln\phi \approx 107$, and the spectral equation (Eq.~3) provides the self-consistency condition linking the lattice scaling factor $\phi$ to the coupling constant $\alpha$.

\subsection{Soliton Creation and Energy Budget (EXP-177)}
\label{sec:soliton_creation}

The DSI potential admits topological solitons at its minima (Section~7.2). We now establish the energy cost and creation mechanism for these solitons on the 600-cell graph.

\textbf{The 5-coloring theorem.} The 5 inscribed 24-cells define a proper 5-coloring of the 600-cell graph: \emph{all} 720 edges connect vertices in different cosets. This is a consequence of the coset structure: if two vertices shared a coset (color), their edge would lie within one of the inscribed 24-cells, but each 24-cell has only 96 edges among its 24 vertices, and no 600-cell edge is an intra-24-cell edge.

\textbf{Result:} $720/720 = 100\%$ of edges are inter-coset. \textbf{Status: DERIVED.}

\textbf{Soliton creation energy.} In the antiferromagnetic 5-state Potts model on the 600-cell graph, flipping a single vertex from its ground-state coset to another coset breaks exactly 3 satisfied edges (the edges to neighbors in the target coset):
\begin{equation}
\boxed{E_{\text{flip}} = 3}
\end{equation}
This energy is \textbf{uniform}: every vertex has exactly 3 neighbors in each of the 4 other cosets, independent of vertex type (A, B, or C).

\textbf{Pair creation.} Creating a soliton--antisoliton pair (two adjacent vertices flipped to each other's cosets) costs:
\begin{equation}
E_{\text{pair}} = 2 \times 3 - 2 = 4
\end{equation}
where the saving of 2 comes from the shared edge between the pair. The pair energy $E_{\text{pair}} = 4$ is also \textbf{uniform} (independent of which edge is chosen).

\textbf{Multi-particle bound states (EXP-165, 166).} The general binding formula for $n$ solitons is:
\begin{equation}
E = 3n - S
\end{equation}
where $S$ counts the edge-savings from adjacency. A triangle of 3 solitons costs $E = 6$ ($S = 3$, binding energy 33\%), and remarkably, a tetrahedron of 4 solitons also costs $E = 4$---the same as a pair. Dense configurations are energetically preferred, consistent with confinement.

\textbf{Charge conservation.} Under the $\mathbb{Z}_5$ coset symmetry, a single soliton flip changes the coset charge, while pair creation conserves it. This provides a topological selection rule for allowed processes.

\subsection{The Soliton Bag: Spectral Properties (EXP-178)}
\label{sec:soliton_bag}

A soliton modifies the local graph structure, creating a ``bag'' region with altered spectral properties.

\textbf{Bag energy.} For a single soliton at vertex $v_0$, the bag is defined as $v_0$ plus its 12 neighbors (13 vertices total). The spectral energy of the bag (sum of Laplacian eigenvalues restricted to the bag) is:
\begin{equation}
\boxed{E_{\text{bag}} = 6.0 \quad \text{(exact, uniform)}}
\end{equation}
This value is independent of vertex type, coset, and target---a remarkable universality given the three different vertex types in the 600-cell.

\textbf{Spectral purity breaking.} The unperturbed 600-cell Laplacian is spectrally pure: all 120 eigenvalues lie in $\mathbb{Q}(\sqrt{5})$. The soliton perturbation partially breaks this purity: 102 of 120 eigenvalues remain in $\mathbb{Q}(\sqrt{5})$, while 18 become irrational (not expressible as $a + b\phi$). The bag's internal Laplacian (on 13 vertices) retains full spectral purity, with eigenvalues $\{0,\, 7-2\phi\;(\times 3),\, 7\;(\times 5),\, 5+2\phi\;(\times 3),\, 13\}$.

\textbf{Localized modes.} The soliton creates 14 modes with participation ratio $\text{PR} < 0.2$ (strongly localized). The highest mode concentrates 55\% of its amplitude at the defect site. This localization provides the mechanism by which solitons acquire mass: the localized spectral weight determines the effective mass through the DSI potential.

\subsection{Cosmological Condensation (EXP-179)}
\label{sec:condensation}

The soliton creation mechanism naturally connects to cosmological particle production through a Kibble-Zurek--type scenario.

\textbf{Critical temperature.} Monte Carlo simulation of the 5-state Potts antiferromagnet on the 600-cell graph reveals a critical temperature $T_c \approx 1$ (in units of the coupling constant), characterized by a sharp transition in soliton density. Below $T_c$, the lattice is fully ordered; above $T_c$, solitons proliferate rapidly.

\textbf{Frozen defects.} Rapid cooling from $T \gg T_c$ to $T = 0$ freezes \textbf{88 topological defects} (out of 120 vertices) in a Kibble-Zurek--type mechanism. These frozen solitons are clustered rather than randomly distributed, reflecting the graph's local connectivity. Slower cooling produces $\sim 88$ defects; faster quenching produces $\sim 95$.

\textbf{Physical interpretation.} In a cosmological context, matter (fermions as coset solitons) is a frozen relic of the hot early universe. As the universe cools through $T_c$, topological defects become trapped at DSI potential minima, acquiring masses $m = m_e \cdot \phi^n$ according to their winding numbers. The pair creation mechanism (Section~\ref{sec:soliton_creation}) ensures that solitons are produced in charge-conserving pairs, and the binding formula $E = 3n - S$ favors the dense multi-particle configurations observed in baryons.

\textbf{Instanton tunneling.} At temperatures below $T_c$, soliton creation proceeds via quantum tunneling (instantons). The Euclidean action is $S_{\text{single}} = 1.31$ for single-soliton creation and $S_{\text{pair}} = 1.75$ for pair creation, giving tunneling rates $\Gamma \propto e^{-S}$ of 0.27 and 0.17 respectively. At $T \approx 3$, the phase-space volume for adjacent pairs ($Z_{\text{pair}} = 190$) exceeds that for single solitons ($Z_{\text{single}} = 141$), making pair production the dominant channel.

\textbf{Status:} The soliton creation energy ($E_{\text{flip}} = 3$), pair energy ($E_{\text{pair}} = 4$), binding formula ($E = 3n - S$), 5-coloring, and $\mathbb{Z}_5$ charge conservation are all \textbf{DERIVED}. The bag energy ($E_{\text{bag}} = 6.0$) and 3-way spectral splitting (Section~\ref{sec:three_splitting}) are \textbf{DERIVED}. The condensation mechanism is \textbf{SPECULATIVE}---the Monte Carlo simulation establishes the existence of frozen defects, but the identification with cosmological particle production requires a full quantum field theory on the 600-cell graph.

\subsection{Connection to Efimov Physics}

The DSI mechanism has a physical analogue in \textbf{Efimov physics}: three-body quantum systems with resonant two-body interactions exhibit a geometric tower of bound states with energy ratios $E_{n+1}/E_n = e^{-2\pi/s_0}$, where $s_0 \approx 1.006$ for identical bosons. The 600-cell lattice generates an analogous structure with $s_0 = \pi/\ln\phi \approx 6.53$ and scaling ratio $\phi \approx 1.618$ (compared to $e^{\pi/1.006} \approx 22.7$ for standard Efimov). The angular frequency is:
\begin{equation}
\omega_{\text{DSI}} = \frac{2\pi}{\ln\phi} \approx 13.057
\end{equation}

The Green's function on the 600-cell graph exhibits oscillatory behavior with a sign change at graph distance $d = 3$ (Section~\ref{sec:interactions}), consistent with log-periodic corrections to long-range potentials. Such corrections are suppressed at macroscopic scales for decay exponents $p \geq 1$.

\textbf{Classification:} The DSI potential (Eq.~\ref{eq:dsi_potential}) and the scaling factor $\phi$ are DERIVED from 600-cell geometry. The mass formula $n = a_1 a + b_1 b$ with $a_1 = 5$, $b_1 = 6$ is DERIVED from intersection numbers. The quantum numbers $(a,b)$ are DERIVED from complexity minimization (Section~\ref{sec:ab_derivation}). The fractal interpretation ($|H_4| = 120^2$) is an OBSERVATION whose physical significance is SPECULATIVE. The Efimov connection and $\omega_{\text{DSI}}$ are DERIVED. The occupied level set $\{0, 3, 5, 11, 16, 17, 19, 26\}$ remains PATTERN.


%==============================================================
\section{E8 Extension and Gauge Unification}
%==============================================================

\subsection{The 600-Cell to E8 Connection via the Icosian Lattice}
\label{sec:e8}

The 240 roots of $E_8$ are constructed explicitly from the 600-cell using the \textbf{icosian lattice} of Conway and Sloane \cite{conway1998icosian}. This construction is intimately related to the Elser-Sloane quasicrystal \cite{elser1987}, where $E_8$ projects to 4D as two concentric 600-cells scaled by $\phi$.

\textbf{Step 1: Decomposition.} Each 600-cell vertex $v \in \mathbb{R}^4$ is decomposed as $v_i = a_i + b_i\phi$ with $a_i, b_i \in \frac{1}{2}\mathbb{Z}$, yielding an 8-tuple $(a_0, b_0, a_1, b_1, a_2, b_2, a_3, b_3) \in \mathbb{R}^8$.

\textbf{Step 2: Icosian norm.} The icosian inner product uses the Gram matrix (per coordinate pair):
\begin{equation}
G = \begin{pmatrix} 1 & 1 \\ 1 & 2 \end{pmatrix}, \qquad \det G = 1 \text{ (unimodular)}
\end{equation}
The icosian norm is $N_{\text{icos}}(v) = \sum_{i=0}^{3}(a_i^2 + 2a_i b_i + 2b_i^2)$. All 120 unit 600-cell vertices have $N_{\text{icos}} = 1$.

\textbf{Step 3: Cholesky embedding.} The Cholesky factorization $G = LL^T$ with $L = \bigl(\begin{smallmatrix}1&0\\1&1\end{smallmatrix}\bigr)$ gives the map $(a_i, b_i) \mapsto (a_i + b_i,\, b_i)$. Under this map, icosian norm equals Euclidean norm$^2$. The 120 vertices embed as a set $S \subset \mathbb{R}^8$ with $\|s\|^2 = 1$.

\textbf{Step 4: Second 600-cell.} The Galois conjugate $\phi' = (1-\sqrt{5})/2$ defines a second copy: $T = \phi' \cdot S$. Since $\phi\phi' = -1$, the conjugation acts on decomposition coefficients as $(a,b) \mapsto (a-b,\,-a)$. The icosian norm of $\phi'$ is $N_{\text{icos}}(\phi') = 1$, so $T$ also consists of unit icosians.

\textbf{Step 5: E8 root system.}
\begin{equation}
\boxed{S \cup T = 240 \text{ vectors in } \mathbb{R}^8, \quad S \cap T = \emptyset}
\end{equation}
The disjointness holds in the 8-dimensional Cholesky embedding. In the original $\mathbb{R}^4$, the coordinate-wise Galois conjugate $\sigma(S)$ overlaps $S$ in exactly 24 vertices (the Galois-fixed points with $b_i = 0$); see below.

After rescaling by $\sqrt{2}$ to match the $E_8$ convention $\|r\|^2 = 2$, the inner product distribution is:
\begin{equation}
\langle r_i, r_j \rangle \in \{-2, -1, 0, +1\}, \quad \text{with multiplicities } 120,\; 6720,\; 15120,\; 6720
\end{equation}
which is the \textbf{exact} $E_8$ root system inner product structure. A Procrustes rotation matches all 240 vectors to the standard $E_8$ basis (240/240 exact).

\textbf{Classification:} DERIVED. The $E_8$ root system is an exact consequence of the 600-cell geometry via the icosian lattice construction. No free parameters.

\textbf{Galois structure and the gauge/fermion split (EXP-156).} The field $\mathbb{Q}(\sqrt{5})$ possesses a Galois automorphism $\sigma: \phi \mapsto \phi' = (1-\sqrt{5})/2$ that acts coordinate-wise on the 600-cell vertices: $v_i = a_i + b_i\phi \;\mapsto\; \sigma(v_i) = a_i + b_i\phi'$. This maps $S$ to a \emph{rotated} copy $\sigma(S)$, also a 600-cell of unit radius on $S^3$ (not a scaled copy---the ``two concentric 600-cells scaled by $\phi$'' of \cite{elser1987} applies to the full infinite $E_8$ lattice, not just the 240~roots).

The Galois-fixed points of $S$ are precisely the 24~vertices with $b_i = 0$ for all~$i$:
\begin{equation}
S \cap \sigma(S) = \{v \in S : b_i = 0\; \forall\, i\} = \text{Type A} \cup \text{Type B} = D_4 \;\text{root system}\; (24\;\text{vertices})
\end{equation}
The 96 Type-C vertices have $b_i \neq 0$ and are \emph{moved} by $\sigma$: they split into 96~$S$-only and 96~$\sigma(S)$-only vertices in $\mathbb{R}^4$, giving $|S \cup \sigma(S)| = 216$ unique points. Each Type-C vertex has a unique Galois partner in $\sigma(S)$ at distance $d \approx 0.270$.

This provides a \textbf{second independent derivation} of the gauge/fermion split: the 24~gauge vertices are the Galois-invariant sector of $\mathbb{Q}(\sqrt{5})$, while the 96~fermion vertices are the Galois-\emph{moving} sector. The algebraic identity $|v|^2 + |\sigma(v)|^2 = 2$ for all vertices, together with $\phi + \phi' = 1$ and $\phi \cdot \phi' = -1$, confirms the arithmetic consistency.

\textbf{Physical interpretation:}
\begin{itemize}
    \item Low energies ($\sim M_Z$): only $S$ (the 600-cell) contributes to the effective geometry
    \item High energies ($\sim M_{\text{GUT}}$): both copies $S \cup T$ become active (full $E_8$)
    \item The Galois conjugation $\phi \to \phi'$ acts as a ``mirror'' symmetry connecting the two copies
    \item The gauge sector ($D_4$) is Galois-invariant; the matter sector is Galois-moving
\end{itemize}

\subsection{The 120-Cell Dual: Spectral Structure and Fibonacci Discriminants}
\label{sec:120cell}

The 120-cell $\{5,3,3\}$ is the dual of the 600-cell $\{3,3,5\}$: its 600 vertices correspond to the 600 tetrahedral cells, with two vertices adjacent when their tetrahedra share a face. A universal \textbf{factor of 3} governs the duality:

\begin{table}[H]
\centering
\begin{tabular}{lccc}
\toprule
\textbf{Property} & \textbf{600-cell} & \textbf{120-cell} & \textbf{Ratio} \\
\midrule
Vertex degree & 12 & 4 & 3 \\
Graph diameter & 5 & 15 & 3 \\
Distinct eigenvalues & 9 & 27 & 3 \\
\bottomrule
\end{tabular}
\caption{The factor-3 duality between the 600-cell and 120-cell.}
\end{table}

The spectral structure of the 120-cell reveals a fundamental distinction. While the 600-cell spectrum is \textbf{spectrally pure}---all eigenvalues lie in $\mathbb{Q}(\sqrt{5})$ and all multiplicities are perfect squares---the 120-cell spectrum involves \emph{five} algebraic number fields:

\begin{table}[H]
\centering
\begin{tabular}{lccl}
\toprule
\textbf{Field} & \textbf{Eigenvalues} & \textbf{Total dim.} & \textbf{Discriminant} \\
\midrule
$\mathbb{Q}$ & 5 & 75 & -- \\
$\mathbb{Q}(\sqrt{5})$ & 10 & 182 & $5 = F_5$ \\
$\mathbb{Q}(\sqrt{2})$ & 2 & 96 & $2 = F_3$ \\
$\mathbb{Q}(\sqrt{13})$ & 2 & 32 & $13 = F_7$ \\
$\mathbb{Q}(\sqrt{21})$ & 2 & 32 & $21 = F_8$ \\
Cubic & 6 & 183 & -- \\
\midrule
Total & 27 & 600 & \\
\bottomrule
\end{tabular}
\caption{Algebraic field decomposition of the 120-cell spectrum. All quadratic discriminants are Fibonacci numbers $F_n$.}
\label{tab:120cell_spectrum}
\end{table}

The \textbf{Fibonacci discriminant pattern} is striking: every quadratic extension appearing in the 120-cell spectrum involves $\sqrt{F_n}$ for Fibonacci numbers $F_3 = 2$, $F_5 = 5$, $F_7 = 13$, $F_8 = 21$. The 600-cell, being spectrally pure in $\mathbb{Q}(\sqrt{5})$, uses only $\sqrt{F_5}$.

Two further structural features emerge:
\begin{enumerate}
\item \textbf{First-six matching:} The six largest eigenvalues of the 120-cell have multiplicities $1^2, 2^2, 3^2, 4^2, 5^2, 6^2$---identical to the 600-cell pattern. The departure occurs only at the 7th eigenvalue.
\item \textbf{Twin cubics:} Two irreducible cubic polynomials $x^3 - x^2 - 7x + 4$ (mult.~$25 = 5^2$) and $x^3 - x^2 - 7x + 8$ (mult.~$36 = 6^2$) share the same $x^2$- and $x$-coefficients, differing only in their constant terms (4 and 8). Galois conjugate eigenvalues always carry the same multiplicity.
\end{enumerate}

The intersection numbers of the 120-cell are $a_1 = 0$, $b_1 = 3$ (cf.\ $a_1 = 5$, $b_1 = 6$ for the 600-cell), and it is \emph{not} fully distance-regular (failing at $d = 4$). This spectral impurity supports the interpretation that the 600-cell is the fundamental geometric object, with the 120-cell playing a derived role.

\subsection{The McKay Correspondence and CKM Exponent Formula (EXP-211--213)}
\label{sec:mckay}

The \textbf{McKay correspondence} provides a deep connection between the binary icosahedral group $2I$ and the $E_8$ root system. Each of the 9 irreducible representations of $2I$ corresponds to a node of the affine $\hat{E}_8$ Dynkin diagram, with the representation dimension matching the Kac label.

\textbf{The 4-way identity.} A remarkable numerical coincidence unifies four independent structures:
\begin{equation}
\boxed{30 = \sum_{j} d_j(2I) = h(E_8) = a_1 \cdot b_1 = |\text{orthogonality classes}|}
\end{equation}
where $\sum_j d_j = 1+2+2+3+3+4+4+5+6 = 30$ is the sum of all irrep dimensions of $2I$, $h(E_8) = 30$ is the Coxeter number, $a_1 \cdot b_1 = 5 \times 6 = 30$ are the 600-cell intersection numbers, and 30 is the number of orthogonality classes in the association scheme.

\textbf{$E_8$ Dynkin legs from McKay branching.} At the $D_4$ branching node (node 4 in Bourbaki numbering, corresponding to the 4-dimensional irrep of $2I$), the affine $\hat{E}_8$ Dynkin diagram splits into three legs:

\begin{table}[H]
\centering
\begin{tabular}{cccl}
\toprule
\textbf{Leg} & \textbf{Nodes (dims)} & \textbf{Leg dim} & \textbf{Geometric identity} \\
\midrule
A (long) & $5+6$ & $11 = a_1 + b_1$ & Diameter + decagons/vertex \\
B (positive) & $3+2+1$ & $6 = b_1$ & Decagons per vertex \\
C (Galois shadow) & $4'+3'+2'$ & $9 = N_{\text{eig}}$ & Number of distinct eigenvalues \\
\bottomrule
\end{tabular}
\caption{$E_8$ Dynkin legs from the McKay correspondence branching at the $D_4$ node. All leg dimensions equal 600-cell invariants.}
\end{table}

\textbf{CKM exponent formula (EXP-205, 213).} The bare CKM mixing exponents $n_{ij}$ satisfy a \emph{unique} linear formula in the generation quantum numbers $(b_1, b_2)$ of the two fermion families:
\begin{equation}
\boxed{n(b_1, b_2) = 9\, b_2 - 5\, b_1 - 6}
\label{eq:ckm_exponent}
\end{equation}
This reproduces all three exponents exactly: $n(0,1) = 9-0-6 = 3$, $n(0,2) = 18-0-6 = 12$, $n(1,2) = 18-5-6 = 7$. The coefficients have a striking interpretation as $E_8$ Dynkin \textbf{leg dimensions}:
\begin{itemize}
    \item $9 = \dim(\text{Leg C})$ (Galois shadow leg)
    \item $5 = a_1 = \dim(\text{Leg A}) - \dim(\text{Leg B}) = 11 - 6$
    \item $6 = \dim(\text{Leg B})$ (positive leg)
\end{itemize}

The Galois automorphism $\sigma: \phi \to \phi'$ swaps Legs B and C. This formula transforms the CKM bare exponents from three independent numbers to a single algebraic structure, with coefficients entirely determined by the $E_8$ McKay geometry.

\textbf{Status:} The 4-way identity is \textbf{DERIVED}. The leg dimensions matching geometric invariants is \textbf{DERIVED}. The CKM exponent formula is \textbf{STRONG PATTERN} (exact numerical match, coefficients explained by $E_8$ legs, but the mechanism selecting this specific linear combination is not yet derived).

\subsection{Unification Results}

Using MSSM beta functions as an $E_8$ proxy:

\begin{table}[H]
\centering
\begin{tabular}{lcc}
\toprule
\textbf{Model} & \textbf{Unification Scale} & \textbf{Discrepancy} \\
\midrule
Standard Model & -- & 27.9\% \\
600-cell + $E_8$ & $\sim 3 \times 10^{17}$ GeV & \textbf{2.4\%} \\
\bottomrule
\end{tabular}
\caption{Gauge unification comparison.}
\end{table}


%==============================================================
\section{Validation and Experimental Predictions}
%==============================================================

\subsection{The Spectral Action}

Following the Spectral Action Principle of Chamseddine and Connes \cite{chamseddine1997}, the total action is defined through the Dirac operator $\mathcal{D}$ on the 600-cell lattice:

\begin{equation}
S = \text{Tr}(f(\mathcal{D}/\Lambda))
\end{equation}

where $\Lambda \sim M_{\text{Planck}}$ is the cutoff scale and $f$ is a cutoff function.

\subsection{Standard Tests (EXP-086)}

\begin{table}[H]
\centering
\begin{tabular}{lcc}
\toprule
\textbf{Test} & \textbf{Result} & \textbf{Error} \\
\midrule
Euler-Lagrange equations & Consistent & -- \\
Lorentz force recovery & Verified & -- \\
Higgs mass & 125.36 GeV & 0.09\% \\
Higgs VEV & 246.2 GeV & 0.0\% \\
Linear dispersion ($n=1$) & \textbf{Excluded} & -- \\
Quadratic dispersion ($n=2$) & Consistent & -- \\
\bottomrule
\end{tabular}
\caption{Standard Lagrangian validation tests.}
\end{table}

\subsection{Advanced Tests (EXP-087)}

\subsubsection{Trans-Planckian Unitarity}

Unitarity is ensured by:
\begin{itemize}
    \item Hermitian Hamiltonian (real Laplacian spectrum)
    \item Finite derivatives on discrete lattice (no Ostrogradsky ghosts)
    \item $H_4$ symmetry (cancellation of negative-norm states)
\end{itemize}

\subsubsection{Weyl Conformal Invariance}

On the 600-cell (approximated as $S^3$):
\begin{itemize}
    \item $\chi = 0 \Rightarrow a \cdot E_4 = 0$ (Euler term vanishes)
    \item Constant curvature $\Rightarrow C_{\mu\nu\rho\sigma} = 0$ (Weyl tensor vanishes)
    \item Conformal anomaly is \textbf{zero} at leading order
\end{itemize}

\subsubsection{Cosmological Constant}

The Heat Kernel expansion suggests a suppression mechanism:

\begin{equation}
\Lambda_{\text{eff}} \sim \Lambda_{\text{natural}} \times \left(\frac{l_P}{L_{\text{Hubble}}}\right)^2 \sim 10^{-122} \times \Lambda_{\text{natural}}
\end{equation}

This yields the correct order of magnitude for the observed cosmological constant.



\subsection{Photon Dispersion from GRBs}

The discrete spacetime structure predicts energy-dependent photon velocities:

\begin{equation}
\Delta t = \xi \left(\frac{E}{E_{\text{Planck}}}\right)^n \frac{d}{c}
\end{equation}

\textbf{Linear model ($n=1$):} \textbf{EXCLUDED}
\begin{itemize}
    \item Prediction: $\Delta t \sim 8$ s for 1 TeV photon at 1 Gpc
    \item Limit from GRB 090510 \cite{fermi2009}: $< 1$ ms
\end{itemize}

\textbf{Quadratic model ($n=2$):} CONSISTENT
\begin{itemize}
    \item Prediction: $\Delta t \sim 10^{-12}$ ms for 1 TeV at 1 Gpc
    \item Testable with CTA (Cherenkov Telescope Array) at higher energies
\end{itemize}

\subsection{Testable Predictions}
\label{sec:testable}

The framework yields over 130 quantitative predictions. We organize the most impactful by experimental accessibility.

\subsubsection{Tier 1: Testable with Existing or Imminent Data}

\begin{enumerate}
    \item \textbf{PMNS $\sin^2\theta_{13} = 1/45$} (error: 0.16\%). The JUNO reactor experiment will measure this quantity with sub-percent precision, providing a direct falsification test. The predicted value 0.02222 lies within $1\sigma$ of the current PDG central value $0.02219 \pm 0.00062$.

    \item \textbf{PMNS $\sin^2\theta_{12} = 2/(\phi+5)$} (error: 0.33\%). JUNO will also refine the solar mixing angle. Predicted: 0.3028.

    \item \textbf{PMNS $\sin^2\theta_{23} = 4/7$} (error: 0.32\%, second octant). DUNE and T2HK will determine the octant and precise value. Predicted: 0.5714.

    \item \textbf{$\Delta m^2_{21}/\Delta m^2_{31} = 4\alpha$} (error: 1.24\%). Relates the neutrino mass-squared ratio to the fine structure constant. JUNO will measure both $\Delta m^2$ values independently.

    \item \textbf{All three CKM angles to $< 0.2\%$} (Appendix~\ref{app:ckm}). Belle~II and LHCb upgrades will refine CKM elements, testing the specific perturbative corrections ($2\alpha_s/3\pi$, $5\alpha_s/\pi$, $\sin^2\theta_W$).
\end{enumerate}

\subsubsection{Tier 2: New Relations Between Observables}

\begin{enumerate}
    \item \textbf{CKM--PMNS factor-3 rule:} $n^{\text{CKM}}_{ij} = 3 \times n^{\text{PMNS}}_{ij}$ for $\theta_{12}$ and $\theta_{13}$, where $3 = N_{\text{colors}}$. Testable as PMNS precision improves.

    \item \textbf{Exact QLC identity:} $\arctan(\phi^{-3}) + \arctan(\phi^{-1}) = \pi/4$ (algebraic identity). Connects the bare Cabibbo angle to the bare PMNS $\theta_{12}$.

    \item \textbf{$m_s = m_\mu$ at $\mu \approx 1.17$~GeV:} The strange quark and muon share the exponent $n=11$, predicting mass equality at a specific QCD scale. Lattice QCD can verify this.

    \item \textbf{CKM exponent formula:} $n_{ij} = 9\,b_j - 5\,b_i - 6$ with coefficients $= E_8$ Dynkin leg dimensions (Section~\ref{sec:mckay}).

    \item \textbf{CP phase:} $\delta_{\text{CKM}} = \arctan(\sqrt{5}) = 65.91^\circ$ (exp: $65.5 \pm 2.5^\circ$, $0.55\%$). The unitarity triangle angle $\gamma = 2\pi/5 = 72^\circ$ (pentagon angle, within $1\sigma$).
\end{enumerate}

\subsubsection{Tier 3: Structural/Falsifiable Predictions}

\begin{enumerate}
    \item \textbf{Zero fourth generation at any energy.} The Generation Theorem (Section~\ref{sec:fibonacci_generations}) proves $N_{\text{gen}} = 3$ from Galois invariance and the Fibonacci sequence. This is \emph{not} an asymptotic freedom bound---it is a theorem. Any collider detection of a 4th-generation fermion would falsify the framework.

    \item \textbf{20 of 27 DSI levels empty.} The level selection rule (Section~\ref{sec:level_selection}) predicts no particles at $n \in \{1,2,4,6,7,8,9,10,12,13,14,15,18,20,21,22,23,24,25,26\}$. Discovery of a particle at any empty level would falsify the framework.

    \item \textbf{Quadratic photon dispersion ($n=2$), not linear.} The discrete spacetime structure predicts $\Delta t \propto (E/E_{\text{Planck}})^2$, with the linear model ($n=1$) already excluded by Fermi GRB data. CTA will probe the quadratic regime.

    \item \textbf{Sterile sector:} The 4-to-3 mechanism predicts a disconnected 4th group of 24 fermion vertices, candidates for right-handed neutrinos or keV-scale dark matter.

    \item \textbf{$U(1)$ topological confinement:} The 48 $U(1)$ edges form a perfect matching with zero Wilson loops and spectral gap $\Delta = 2$ (Section~\ref{sec:u1_confinement}).
\end{enumerate}

%==============================================================
\section{Summary, Limitations, and Conclusions}
\label{sec:conclusions}
%==============================================================



\begin{table}[H]
\centering
\begin{tabular}{lccc}
\toprule
\textbf{Quantity} & \textbf{Formula} & \textbf{Error} & \textbf{Status} \\
\midrule
\multicolumn{4}{l}{\textit{Coupling constants:}} \\
$\alpha$ & Spectral equation & 0.0001\% & Derived \\
$\alpha_s$ & $1/(2\phi^3)$ & 0.03\% & Derived \\
$\alpha_s/\alpha$ & $10\phi$ & Exact & Derived \\
$\sin^2\theta_W$ & $6/26$ & 0.19\% & Derived \\
$m_H$ & $m_W(\phi - 8\alpha)$ & 0.09\% & Derived \\
$\alpha_G$ & $\alpha^{8\phi^2}$ & 0.5\% & Derived \\
$m_e/m_P$ & $\alpha^{4\phi^2}$ & 0.16\% & Derived \\
\midrule
\multicolumn{4}{l}{\textit{Lepton mass hierarchy (from 600-cell holonomy):}} \\
$m_\mu/m_e$ & $\phi^{11}$ ($n=5+6$) & 3.8\% & \textbf{Derived} \\
$m_\mu/m_e$ (corrected) & $\phi^{5+\sin^2\theta_W+6-\phi^{-4}}$ & \textbf{0.28\%} & \textbf{Derived} \\
$m_\tau/m_\mu$ & $\phi^{6}$ & 6.7\% & \textbf{Derived} \\
$m_\tau/m_e$ (corrected) & $\phi^{d_{\text{eff}}+2k_{\text{eff}}}$ & \textbf{0.25\%} & \textbf{Derived} \\
$m_\tau/m_e$ & $\phi^{17}$ ($n=11+6$) & 2.7\% & \textbf{Derived} \\
$m_u/m_e$ & $\phi^{3}$ & 0.2\% & \textbf{Derived} \\
\midrule
\multicolumn{4}{l}{\textit{Unified fermion mass formula:}} \\
Full formula & $m_f = m_e \cdot \phi^{a \cdot d_{\text{eff}} + b \cdot k_{\text{eff}}}$ & 2.3\% avg & Derived \\
$m_c$ (charm) & $(2,1)$, up-type corr. & 1.03\% & Derived \\
$m_b$ (bottom) & $(-1,4)$, down-type corr. & 0.26\% & Derived \\
$m_t$ (top) & $(4,1)$, up-type corr. & 1.82\% & Derived \\
\midrule
\multicolumn{4}{l}{\textit{DSI mechanism (Section~\ref{sec:dsi}):}} \\
Scaling factor & $\phi$ from edge angle $\pi/5$ & Exact & Derived \\
DSI potential & $V \propto \sin^2(\pi\ln|\psi|/\ln\phi)$ & -- & Derived \\
Level selection & $n = 5a + 6b$ & -- & Derived \\
$(a,b)$ assignment & $\arg\min|a|+|b|$ s.t. $5a+6b=n$ & 9/9 & Derived \\
\midrule
\multicolumn{4}{l}{\textit{E8 and level selection (Sections~\ref{sec:e8}, \ref{sec:level_selection}):}} \\
$E_8$ roots & $S \cup T$ (icosian) & 240/240 & \textbf{Derived} \\
Gauge group & 24-cell $\to D_4 \xrightarrow{E_8}$ SM (Kac $\{3,4,5\}$) & Exact & \textbf{Derived} \\
Level selection & $|N| \in \{0,1,5,19\}$ + 3 lines + $N>0$ & 31/31$^*$ & Pattern+Derived \\
\midrule
\multicolumn{4}{l}{\textit{Interaction structure (Section~\ref{sec:interactions}):}} \\
Selection rules & A-A $=$ A-B $=$ B-B $= 0$ & Exact & \textbf{Derived} \\
Triangle types & ACC : BCC : CCC $= 1:2:2$ & Exact & \textbf{Derived} \\
Per-fermion ratio & $5:10:15 = 1:2:3$ & Exact & \textbf{Derived} \\
Propagator sign & Oscillatory at $d \geq 3$ & -- & Derived \\
Edge split & $1 + 3 + 8 = \dim(\text{SM})$ & Exact & \textbf{Derived} \\
Eff.\ gauge adj. & $A_{\text{eff}} = A_{\text{24-cell}} = A^2|_g/3$ & Exact & \textbf{Derived} \\
Gram spectrum & $\{36,24,12,6,0\} \times \{1,4,9,8,2\}$ & Exact & \textbf{Derived} \\
Pure plaquettes & Only $SU(3)$: 288 of 1200 & Exact & \textbf{Derived} \\
Separation thm. & Form (topology) $\perp$ values (spectrum) & -- & \textbf{Derived} \\
Shared-nbr profiles & Exact per edge type (std $= 0$) & Exact & \textbf{Derived} \\
$10\phi$ decomp. & $480/48 \cdot \phi = 10\phi$ & Exact & \textbf{Derived} \\
$SU(3)$ sub-types & T3: 8 comp.\ $= |A| = \dim(SU(3))$ & Exact & \textbf{Derived} \\
Non-commutativity & $[A_{T_i}, A_{T_j}] \neq 0$ & Exact & \textbf{Derived} \\
A-vertex mechanism & 8 T3 comp.\ $=$ 8 A-vertices & Exact & \textbf{Derived} \\
Mixed-scale ratio & $\alpha_s(M_Z)/\alpha(0) = 10\phi$ & 0.03\% & Clarifying \\
$U(1)$ confinement & Perfect matching, gap $= 2$ & Exact & \textbf{Derived} \\
Gauge hierarchy & Range: $U(1) < SU(2) < SU(3)$ & -- & \textbf{Derived} \\
\midrule
\multicolumn{4}{l}{\textit{Soliton physics (Sections~\ref{sec:soliton_creation}--\ref{sec:condensation}):}} \\
5-coloring & 720/720 edges inter-coset & Exact & \textbf{Derived} \\
$E_{\text{flip}}$ & $3$ (uniform, all vertices) & Exact & \textbf{Derived} \\
$E_{\text{pair}}$ & $4 = 2 \times 3 - 2$ & Exact & \textbf{Derived} \\
Binding formula & $E = 3n - S$ & Exact & \textbf{Derived} \\
$E_{\text{bag}}$ & $6.0$ (uniform, all types) & Exact & \textbf{Derived} \\
3-way splitting & $(m_k-3,\, 2,\, 1)$ universal & Exact & \textbf{Derived} \\
Generation Theorem & $N_{\text{gen}} = 3$ (Galois + Fibonacci) & Exact & \textbf{Derived} \\
\midrule
\multicolumn{4}{l}{\textit{Mixing angles:}} \\
$\theta_{12}$ (CKM) & $\arctan(\phi^{-3}(1 - 2\alpha_s/3\pi))$ & \textbf{0.001\%} & Pattern \\
$\theta_{23}$ (CKM) & $\arctan(\phi^{-7}(1 + 5\alpha_s/\pi))$ & \textbf{0.17\%} & Pattern \\
$\theta_{13}$ (CKM) & $\arctan(\phi^{-12}(1 + \sin^2\theta_W))$ & \textbf{0.096\%} & Pattern \\
$\delta_{\text{CKM}}$ & $\arctan(\sqrt{5})$ & 0.55\% & Pattern \\
CKM--PMNS factor & $n^{\text{CKM}} = 3 \times n^{\text{PMNS}}$ & Exact$^\dagger$ & Pattern \\
$\sin^2\theta_{13}$ (PMNS) & $1/45 = 1/(d \times N_{\text{eig}})$ & \textbf{0.16\%} & Pattern \\
$\sin^2\theta_{12}$ (PMNS) & $2/(\phi+5)$ & \textbf{0.33\%} & Pattern \\
$\sin^2\theta_{23}$ (PMNS) & $4/7$ & \textbf{0.32\%} & Pattern \\
$\Delta m^2$ ratio & $\Delta m^2_{21}/\Delta m^2_{31} = 4\alpha$ & 1.24\% & Pattern \\
\midrule
\multicolumn{4}{l}{\textit{McKay correspondence and branching (Section~\ref{sec:mckay}, \ref{sec:branching_mixing}):}} \\
4-way identity & $30 = \Sigma d_j = h(E_8) = a_1 b_1$ & Exact & \textbf{Derived} \\
CKM exponents & $n = 9b_2 - 5b_1 - 6$ (leg dims) & Exact & Strong Pattern \\
Bare mixing & $p = 1/6$ (Schur's lemma) & Exact & \textbf{Derived} \\
$\rho_8$ branching & $\psi_3 \oplus \psi_4 \oplus \psi_5$ & Exact & \textbf{Derived} \\
\midrule
\multicolumn{4}{l}{\textit{Mass correction structure (Section~\ref{sec:galois_complementarity}):}} \\
Lepton Galois dir. & $\delta_d/\delta_k \approx -\phi$ & 1.7\% & Pattern \\
Up quark $\delta_k$ & $\approx \alpha_s$ & 0.5\% & Pattern \\
Galois compl. & Lepton: $\phi'$ only; up: $\phi$ only & -- & Pattern \\
Galois flip & $\delta z'(\text{lep})/\delta z(\text{up}) \approx 1$ & 0.55\% & Pattern \\
$A$ (Wolfenstein) & $\phi^{-1/2}$ & 0.5\% & Pattern \\
$\gamma_{\text{UT}}$ & $2\pi/5$ (pentagon) & $<1\sigma$ & Pattern \\
$n_{ij}$ from $a_1, b_1$ & $\{3,7,12\}$ & Exact & Strong Pattern \\
Bare QLC & $\arctan(\phi^{-3})+\arctan(\phi^{-1})=\pi/4$ & Exact & \textbf{Derived} \\
\bottomrule
\end{tabular}
\caption{Summary of results. $^*$31/31 with per-line refinement; the $b \leq 2$ constraint (three generations) is now derived via Galois invariance (Section~\ref{sec:fibonacci_generations}). $^\dagger$Exact for $\theta_{12}$ and $\theta_{13}$; $\theta_{23}$ is anomalous.}
\end{table}



The 600-cell framework provides a coherent geometric basis for understanding the origin of fundamental constants in the particle physics sector. The key achievements are:

\begin{enumerate}
    \item \textbf{Numerical precision:} The coupling constants $\alpha$ (0.0001\%), $\alpha_s$ (0.03\%), $\sin^2\theta_W$ (0.19\%), and the Higgs mass (0.09\%) are derived from spectral and geometric properties of the 600-cell, not fitted.

    \item \textbf{Lepton mass hierarchy:} Exponents $n=11=5+6$ and $n=17=11+6$ are \textbf{derived} from the graph diameter (5) and decagons per vertex (6)---topological invariants of the 600-cell. Electroweak phase corrections ($\delta_d = \sin^2\theta_W$, $\delta_k = -1/\phi^4$) reduce errors to \textbf{0.25--0.53\%}.

    \item \textbf{Unified mass formula:} $m_f = m_e \cdot \phi^{a \cdot d_{\text{eff}} + b \cdot k_{\text{eff}}}$ predicts seven of nine fermion masses within 2\%. The sector-dependent corrections use only quantities already derived within the theory.

    \item \textbf{DSI mechanism:} The mass formula is not an empirical fit but arises from a Discrete Scale Invariance potential (Eq.~\ref{eq:dsi_potential}) whose scaling factor $\phi$ is fixed by the 600-cell geometry. Fermions are topological solitons whose mass levels are selected by their winding numbers $(a,b)$ on the lattice. Crucially, $(a,b)$ are not free parameters: they are uniquely determined by complexity minimization on the graph (Section~\ref{sec:ab_derivation}).

    \item \textbf{Spectral structure:} The 600-cell adjacency spectrum has 9 distinct eigenvalues, all of the form $a + b\phi$, with multiplicities that are all perfect squares $n^2$ ($n = 1, \ldots, 6$). The intersection numbers $a_1 = 5$, $b_1 = 6$ provide a graph-theoretic derivation of $\alpha_s/\alpha = 10\phi$ and a second independent path to $\sin^2\theta_W = 6/26$.

    \item \textbf{Gauge group derivation (complete):} The 24 non-fermionic vertices form the $D_4$ root system. The $E_8$ embedding breaks $D_4$ triality: the three legs extend with tail lengths $\{0, 1, 3\}$ and distinct Kac labels $\{3, 4, 5\}$, uniquely assigning each leg to a SM factor (Section~\ref{sec:gauge}). The final step---$SU(2)_R \to U(1)_Y$---follows from the dead-end isolation of the $m=3$ leg: with no non-Abelian extension to protect it, $SU(2)_R$ naturally reduces to its Cartan subalgebra $U(1)_Y$. Three independent structural arguments (topological isolation, maximal commutant $A_7$ with dim~63, smallest Kac label) confirm the breaking. The full chain is: 600-cell $\to$ 24-cell $\to$ $D_4$ $\to$ $E_8$ Dynkin $\to$ SM.

    \item \textbf{Fermion mapping:} Three generations encoded in $96 = 16 \times 3 \times 2$ vertices, with a geometric 4-to-3 selection mechanism via the temporal axis.

    \item \textbf{Mixing angles:} CKM exponents $\{3, 7, 12\}$ satisfy the unique linear formula $n = 9b_2 - 5b_1 - 6$, with coefficients equal to $E_8$ Dynkin leg dimensions (Section~\ref{sec:mckay}). Cross-pairings of $H_4$ and $H_4'$ exponents from the $E_8$ branching provide a cleaner structure for the half-integer refinement (Appendix~\ref{app:ckm}). The CP-violating phase $\delta_{\text{CKM}} = \arctan(\sqrt{5}) = 65.91^\circ$ matches experiment to $0.55\%$ ($< 1\sigma$). A factor-3 rule $n^{\text{CKM}} = 3 \times n^{\text{PMNS}}$ holds for $\theta_{12}$ and $\theta_{13}$ (Appendix~\ref{app:ckm}). All three PMNS mixing angles are predicted by purely geometric $\sin^2$ formulas: $\sin^2\theta_{13} = 1/45$ ($0.16\%$), $\sin^2\theta_{12} = 2/(\phi+5)$ ($0.33\%$), $\sin^2\theta_{23} = 4/7$ ($0.32\%$)---all within $1\sigma$ of PDG~2024 (Table~\ref{tab:pmns_geometric}).

    \item \textbf{McKay correspondence (new):} The 4-way identity $30 = \sum d_j(2I) = h(E_8) = a_1 \cdot b_1$ unifies representation theory, Coxeter theory, and 600-cell graph invariants. The three $E_8$ Dynkin legs at the $D_4$ branching node have dimensions $(11, 6, 9) = (a_1 + b_1,\, b_1,\, N_{\text{eig}})$, all equal to 600-cell invariants (Section~\ref{sec:mckay}).

    \item \textbf{Bare flavor mixing (new):} The branching $\rho_8|_{2T} = \psi_3 \oplus \psi_4 \oplus \psi_5$ gives three generations via three 2D irreps (two complex conjugate, one real). The bare mixing matrix is democratic with $p = 1/6$, derived from Schur's lemma and the 2-transitivity of $2I$ on cosets (Section~\ref{sec:branching_mixing}). Physical CKM/PMNS require mass-induced symmetry breaking.

    \item \textbf{120-cell dual:} The dual polytope's spectrum spans five algebraic fields, with all quadratic discriminants being Fibonacci numbers ($F_3, F_5, F_7, F_8$). This spectral impurity confirms the 600-cell as the fundamental object.

    \item \textbf{E8 from 600-cell:} The icosian lattice construction (Section~\ref{sec:e8}) produces the 240 $E_8$ roots exactly from two copies of the 600-cell ($S$ and $T = \phi' \cdot S$), with zero overlap and exact inner product matching. This is a derivation, not a coincidence.

    \item \textbf{Level selection and three generations:} Two algebraic conditions---$\mathbb{Z}[\phi]$ norm restriction and a three-line rule---plus norm positivity achieve \textbf{31/31} classification accuracy (Section~\ref{sec:level_selection}). The three-generation limit ($b \leq 2$ on $a=1$) is now \textbf{derived}: Galois invariance of the $E_8$ lattice forces a dual-sector energy functional $E = V(z) + V(z')$, whose vanishing requires $z$ to be a unit of $\mathbb{Z}[\phi]$; Dirichlet's unit theorem and the Fibonacci condition $F(k\!-\!1)=1$ then give exactly $b \in \{0,1,2\}$ (Section~\ref{sec:fibonacci_generations}).

    \item \textbf{Gauge unification:} $E_8$ extension reduces discrepancy from 27.9\% to 2.4\%.

    \item \textbf{Validated Lagrangian:} Passes all tests---unitarity, Weyl invariance, correct Higgs sector.

    \item \textbf{Interaction selection rules:} The edge structure automatically forbids A-A, A-B, and B-B connections, leaving only three types of interaction vertex (ACC, BCC, CCC) in the exact ratio $1:2:2$. Each fermion participates in $5 + 10 + 15 = 30$ triangles with ratio $1:2:3$. The massless propagator changes sign at graph distance $d = 3$, exhibiting confinement-like behavior on the finite lattice (Section~\ref{sec:interactions}). The $U(1)$ sector is topologically confined: its 48 edges form a perfect matching with zero cycle rank and maximum spectral gap $\Delta = 2$ (Section~\ref{sec:u1_confinement}). The three gauge channels exhibit a derived hierarchy: $U(1)$ (range 1) $< SU(2)$ (range 2) $< SU(3)$ (long-range), matching the physical expectation (Section~\ref{sec:gauge_hierarchy}).

    \item \textbf{Edge split $\mathbf{1+3+8}$ (new):} Among the 9 C-neighbors of each fermion, exactly one has a geometrically unique gauge-overlap profile (degree 4 in local CC graph, shared A but zero shared B). Reclassifying this single edge yields $1\,(\text{AC}) + 3\,(2\,\text{BC} + 1\,\text{CC}_*) + 8\,\text{CC} = 12 = \dim(\text{SM})$. The effective gauge adjacency through shared fermions reproduces the 24-cell \emph{exactly}: $A_{\text{eff}} = A^2|_{\text{gauge}}/3$ (Section~\ref{sec:edge_split}).

    \item \textbf{Plaquette purity and separation theorem (new):} Among 1200 triangular plaquettes, only $SU(3)$ has pure plaquettes (288)---$U(1)$ and $SU(2)$ have zero---reproducing the non-Abelian self-interaction structure of QCD. A vertex-transitivity theorem proves that spectral projections are sector-blind ($\chi_k$ identical across U(1), SU(2), SU(3) for all 9 eigenspaces), establishing that the Lagrangian \emph{form} (from topology) and coupling \emph{values} (from the spectrum) are two independent geometric properties of the 600-cell (Section~\ref{sec:edge_split}).

    \item \textbf{Shared-neighbor asymmetry and topological $10\phi$ (new):} Each edge type has an \emph{exact} shared-neighbor profile (standard deviation zero): $U(1)$ and $SU(2)_{\text{ch}}$ edges share only fermion (C-type) neighbors, while $SU(3)$ edges share 1.75 gauge-boson neighbors on average---the topological origin of Abelian vs.\ non-Abelian character. The ratio $\alpha_s/\alpha = 10\phi$ decomposes as $(\text{fermion-shared}_{U(1)} / \text{gauge-shared}_{SU(2)}) \cdot \phi = (480/48) \cdot \phi$, giving the factor 10 a topological interpretation independent of the spectral derivation (Section~\ref{sec:edge_split}).

    \item \textbf{$SU(3)$ sub-type decomposition and A-vertex mechanism (new):} The 384 $SU(3)$ edges split into three sub-types by shared-neighbor profile: $T_1(0A,1B,4C) = 96$ edges with 24 components, $T_2(0A,2B,3C) = 96$ edges with 32 components, and $T_3(1A,1B,3C) = 192$ edges with \textbf{8 connected components}. The mechanism is now \emph{derived}: each T3 edge shares exactly one Type-A vertex, and the 8 A-vertices partition T3 into 8 disjoint connected groups of 24 edges, giving $\dim(SU(3)) = |A| = 8$ as a geometric identity (EXP-149). The sub-type adjacency matrices do not commute ($[A_{T_i}, A_{T_j}] \neq 0$), reproducing non-Abelian structure (Section~\ref{sec:edge_split}).

    \item \textbf{Soliton creation mechanism (new):} The 600-cell graph admits a proper 5-coloring (all 720 edges inter-coset), and coset solitons have a uniform creation energy $E_{\text{flip}} = 3$. Pair creation costs $E_{\text{pair}} = 4$ with $\mathbb{Z}_5$ charge conservation. The general binding formula $E = 3n - S$ explains why dense multi-particle states (baryons) are energetically preferred (Section~\ref{sec:soliton_creation}).

    \item \textbf{Soliton bag model (new):} The spectral energy of the 13-vertex soliton bag is $E_{\text{bag}} = 6.0$ exactly, uniform across all vertex types and cosets. The soliton partially breaks the spectral purity of the 600-cell (102/120 eigenvalues remain in $\mathbb{Q}(\sqrt{5})$), while the bag's internal spectrum retains full purity (Section~\ref{sec:soliton_bag}).

    \item \textbf{Spectral derivation of three generations (new):} A single soliton defect (rank-3 perturbation) splits each eigenvalue group into exactly three sub-levels with universal pattern $(m_k - 3, 2, 1)$. The three sub-levels have distinct localization: bulk (delocalized), shell (neighbor-only, $|\psi(v_0)|^2 = 0$ exactly), and core (defect-localized). This provides a second independent derivation of $N_{\text{gen}} = 3$ from the uniform coset structure $3 = 12/4$ (Section~\ref{sec:three_splitting}).

    \item \textbf{Generation Theorem $N_{\text{gen}} = 3$ (new):} The number of generations is \textbf{derived} in five steps: (1)~600-cell spectrum places levels in $\mathbb{Z}[\phi]$; (2)~$E_8$ icosian construction provides a Galois shadow $z' = \sigma(z)$; (3)~Galois invariance of $E_8$ (defined over $\mathbb{Z}$) forces the energy $E = V(z) + V(z')$, whose vanishing implies $|N(z)| = 1$ (unit); (4)~Dirichlet's theorem gives units $= \pm\phi^k$, and $\phi^k$ on $a=1$ requires $F(k\!-\!1) = 1$ (three Fibonacci solutions); (5)~$b \in \{0,1,2\}$. The triple identity $a_1 = 5 = \operatorname{disc}(\mathbb{Z}[\phi]) = |2I/2T|$ unifies graph, algebra, and group theory (Section~\ref{sec:fibonacci_generations}).
\end{enumerate}

\textbf{Limitations and open problems:}
\begin{itemize}
    \item The number $N_{\text{gen}} = 3$ is now derived via three independent routes: (i) Galois invariance $+$ Fibonacci (Section~\ref{sec:fibonacci_generations}), (ii) soliton spectral splitting (Section~\ref{sec:three_splitting}), and (iii) the 4-to-3 temporal axis mechanism. A fourth route via representation theory ($\rho_8|_{2T} = \psi_3 \oplus \psi_4 \oplus \psi_5$, three 2D irreps) reinforces this from McKay/group theory (Section~\ref{sec:branching_mixing}). However, connecting the three sub-levels to the three \emph{specific mass scales} (electron/muon/tau families) remains at the PATTERN level: the localization hierarchy (bulk/shell/core) qualitatively matches the Yukawa hierarchy, but quantitative mass ratios are not reproduced.
    \item \textbf{Bare mixing is democratic:} The group-theoretic mixing matrix $M = (1/6)I + (5/12)(J-I)$ is fully $S_3$-symmetric (Section~\ref{sec:branching_mixing}). Deriving the physical CKM and PMNS matrices requires identifying the symmetry-breaking mechanism that lifts this democracy. The structural asymmetry $\psi_3, \psi_4$ (complex) vs.\ $\psi_5$ (real) provides a candidate distinction but does not yet yield quantitative mixing angles.
    \item With perturbative gauge corrections, all three CKM mixing angles are now predicted to $< 0.2\%$ (Appendix~\ref{app:ckm}, Table~\ref{tab:ckm_corrected}). The bare exponents $\{3, 7, 12\}$ are expressible through $H_4$ Coxeter exponents but not uniquely determined. The correction coefficients---$2\alpha_s/(3\pi)$ for $\theta_{12}$, $5\alpha_s/\pi$ for $\theta_{23}$, and $\sin^2\theta_W$ for $\theta_{13}$---involve only quantities derived within this framework, but the specific association of each correction to each angle is not yet derived. The CP phase $\delta_{\text{CKM}} = \arctan(\sqrt{5})$ and the factor-3 rule $n^{\text{CKM}} = 3 \times n^{\text{PMNS}}$ are strong patterns but not yet derived. PMNS exponents ($n = 0, 1, 4$) are motivated by geometric dimensions.
    \item The gauge group derivation is now complete (EXP-135): $D_4$ from 24-cell, triality broken by $E_8$ Dynkin topology, leg assignment via Kac labels $\{3,4,5\}$, and $SU(2)_R \to U(1)_Y$ from dead-end isolation. The remaining open questions are the breaking \emph{scale} (the $W_R$ mass) and hypercharge normalization.
    \item \textbf{Neutrino masses} are not addressed by the current framework. The second 600-cell $T = \phi' \cdot S$ in the $E_8$ construction does not produce viable neutrino masses: the Galois identity $\phi \cdot \phi' = -1$ kills level-dependent mass formulas, and Type-I seesaw with $m_D = m_{\text{lepton}}$ gives a hierarchy too steep by a factor $\sim 2450$ (EXP-134). The neutrino sector likely requires a different mechanism (Majorana masses, radiative generation, or the $(1,8)$ sector of $E_8$ branching).
    \item The Down quark mass prediction (5.15\,MeV) lies at the 1$\sigma$ upper bound of the PDG value ($4.67^{+0.48}_{-0.17}$\,MeV), placing it within $1\sigma$ of experiment (EXP-128). The Down--Strange tension is structural and irreducible within shared corrections.
    \item The $E_8$ unification uses MSSM beta functions as a proxy; a direct derivation from $E_8 \to H_4$ breaking is needed.
    \item \textbf{Gravity} is not addressed by this framework; the 600-cell encodes particle physics structure but does not produce a viable gravitational theory.
    \item \textbf{Coupling constant scales (EXP-150):} The ratio $\alpha_s/\alpha = 10\phi$ uses mixed renormalization scales ($\alpha$ at $q = 0$ Thomson limit, $\alpha_s$ at $M_Z$). At a single scale $M_Z$, the ratio is 15.09 (6.8\% off from $10\phi$). The framework encodes coupling constants at their natural measurement scales, consistent with a non-perturbative geometric origin rather than perturbative running.
    \item \textbf{Electroweak scale (EXP-152):} The framework derives mass \emph{ratios} ($m_H/m_W$, $m_Z/m_W$) but not the absolute electroweak scale $v = 246\,$GeV, which remains an additional input. The top Yukawa coupling $y_t = 0.992 \approx 1$ and $n_{\text{top}} = 26 = \dim(\text{phi-sector})$ suggest a special role for the top quark, but no derivation of $v$ from the 600-cell has been found.
\end{itemize}

Despite these limitations, the framework offers a concrete, falsifiable geometric hypothesis for the origin of fundamental constants. The derivation of $N_{\text{gen}} = 3$ from Galois invariance and the Fibonacci sequence---together with the coupling constant precision ($\alpha$ at $0.0001\%$, $\alpha_s$ at $0.03\%$), the lepton mass hierarchy ($0.25\%$ after corrections), all CKM angles below $0.2\%$, and all PMNS angles within $1\sigma$ of PDG~2024---suggest that the 600-cell geometry captures genuine structure of physical law, even if the complete picture remains incomplete. The McKay correspondence connecting $2I$ irreps to the $E_8$ Dynkin diagram, and the resulting CKM exponent formula $n = 9b_2 - 5b_1 - 6$ with $E_8$ leg dimensions as coefficients, provide a new algebraic link between the generation structure and mixing angles.

%==============================================================
\section*{Acknowledgments}
%==============================================================

We acknowledge the foundational work on spectral geometry by Alain Connes and Ali Chamseddine \cite{chamseddine1997}, the geometric optimality results of Cohn and Kumar \cite{cohn2007,cohn2011}, and the 600-cell GUT framework developed by Christopher C. O'Neill \cite{oneill2025gut,oneill2025120cell,oneill2023binary}. We also acknowledge the experimental data from CODATA 2022 \cite{codata2022}, the Particle Data Group \cite{pdg2024}, and the Fermi LAT collaboration \cite{fermi2009}.

%==============================================================
\appendix
%==============================================================

%==============================================================
\section{CKM and PMNS Mixing Matrices}
\label{app:ckm}
%==============================================================

\textbf{Note:} The mixing angle predictions presented here are at the PATTERN level---the exponents connect to $H_4$ Coxeter exponents and geometric dimensions but are not uniquely derived from the 600-cell geometry.


\subsection{The Phi-Tower Formula}

Mixing angles follow a geometric hierarchy based on powers of the golden ratio:

\begin{equation}
\boxed{\theta_{ij} = \arctan(\phi^{-n_{ij}})}
\end{equation}

\subsection{CKM Matrix (Quarks)}

\textbf{Integer exponents.} The leading-order CKM hierarchy follows the phi-tower formula with integer exponents motivated by algebraic dimensions:

\begin{table}[H]
\centering
\begin{tabular}{ccccc}
\toprule
\textbf{Angle} & \textbf{$n$} & \textbf{Calculated} & \textbf{Experimental} & \textbf{Error} \\
\midrule
$\theta_{12}$ (Cabibbo) & 3 & $13.28°$ & $13.04°$ & \textbf{1.8\%} \\
$\theta_{23}$ & 7 & $1.97°$ & $2.34°$ & 15.6\% \\
$\theta_{13}$ & 12 & $0.18°$ & $0.22°$ & 18.7\% \\
\bottomrule
\end{tabular}
\caption{CKM mixing angles from phi-tower with integer exponents.}
\end{table}

\textbf{Coxeter exponent origin.} The exponents $\{3, 7, 12\}$ are expressible through the $H_4$ Coxeter exponents $\{1, 11, 19, 29\}$:
\begin{equation}
3 = \frac{1+11}{4}, \qquad 7 = \frac{|1-29|}{4}, \qquad 12 = \frac{19+29}{4}
\end{equation}
These are also compatible with algebraic dimension assignments: $n = 3 = \dim(\text{Im}(\mathbb{H}))$, $n = 7 = \dim(\text{Im}(\mathbb{O}))$, $n = 12 = $ vertex degree.

\textbf{Half-integer refinement.} Allowing half-integer exponents, a scan yields $\sin\theta_{ij} = \phi^{-n_{ij}}$ with $n = (3,\, 13/2,\, 23/2)$, giving mean error $\textbf{6.0\%}$ across all 9 CKM elements. The bare geometric prediction $\theta_C = \arctan(\phi^{-3}) = 13.28°$ has $1.8\%$ error. Including the standard 1-loop QCD vertex correction for quarks in the fundamental representation:
\begin{equation}
\boxed{\theta_C = \arctan\!\left(\phi^{-3}\left(1 - \frac{2\alpha_s}{3\pi}\right)\right) = 12.962°}
\label{eq:cabibbo_qcd}
\end{equation}
matching the experimental value $12.96°$ to $\textbf{0.001\%}$ (EXP-191). The factor $2\alpha_s/(3\pi)$ is exactly the QCD correction to the quark--$W$ vertex at one loop, providing a physical interpretation: the bare angle is geometric ($\phi^{-3}$), while the QCD dressing shifts it to the observed value.

\textbf{Perturbative corrections to all three CKM angles (EXP-192).} The same philosophy---geometric bare value plus perturbative gauge correction---extends to the remaining two angles:
\begin{align}
\theta_{23} &= \arctan\!\left(\phi^{-7}\left(1 + \frac{5\alpha_s}{\pi}\right)\right) = 2.342° && \text{(exp: } 2.338°\text{, error \textbf{0.17\%})} \label{eq:theta23_qcd} \\
\theta_{13} &= \arctan\!\left(\phi^{-12}\left(1 + \sin^2\theta_W\right)\right) = 0.2191° && \text{(exp: } 0.2189°\text{, error \textbf{0.096\%})} \label{eq:theta13_ew}
\end{align}
The correction to $\theta_{23}$ is $5\alpha_s/\pi$, where the coefficient 5 equals both the first intersection number $a_1$ and the discriminant of $\mathbb{Z}[\phi]$. The correction to $\theta_{13}$ is $\sin^2\theta_W = 6/26$, which is already derived within this framework (Section~\ref{sec:spectral}). These results transform the CKM sector from a single good prediction ($\theta_{12}$, $1.8\%$) to three high-precision predictions (all $< 0.2\%$):

\begin{table}[H]
\centering
\begin{tabular}{lcccc}
\toprule
\textbf{Angle} & \textbf{Bare (error)} & \textbf{Correction} & \textbf{Corrected (error)} & \textbf{Type} \\
\midrule
$\theta_{12}$ & $\phi^{-3}$ (1.8\%) & $1 - 2\alpha_s/(3\pi)$ & $12.962°$ (\textbf{0.001\%}) & QCD 1-loop \\
$\theta_{23}$ & $\phi^{-7}$ (15.6\%) & $1 + 5\alpha_s/\pi$ & $2.342°$ (\textbf{0.17\%}) & QCD \\
$\theta_{13}$ & $\phi^{-12}$ (18.7\%) & $1 + \sin^2\theta_W$ & $0.2191°$ (\textbf{0.096\%}) & Electroweak \\
\bottomrule
\end{tabular}
\caption{CKM mixing angles: bare geometric values with perturbative corrections. The corrections involve only quantities already derived in this framework ($\alpha_s = 1/2\phi^3$, $\sin^2\theta_W = 6/26$).}
\label{tab:ckm_corrected}
\end{table}

\textbf{Status:} The corrected angles are classified as PATTERN. The bare exponents $\{3, 7, 12\}$ are not yet derived from first principles, though they match the Coxeter structure below. The correction coefficients (particularly the factor 5 in $\theta_{23}$) require further investigation.

\textbf{$E_8$ cross-pairing (EXP-129).} The half-integer exponents also arise from \emph{cross-pairings} of $H_4$ and $H_4'$ Coxeter exponents:
\begin{equation}
n_{12} = \frac{1+11}{4} = 3, \qquad n_{23} = \frac{19+7}{4} = 6.50, \qquad n_{13} = \frac{29+17}{4} = 11.50
\end{equation}
where $\{1,11,19,29\}$ are $H_4$ exponents and $\{7,17\}$ are $H_4'$ exponents from the $E_8$ branching $E_8 \to H_4 \oplus H_4'$. The cross-pairing $n_{23} = 6.50$ matches the fitted value $6.65$ to $2.3\%$, and $n_{13} = 11.50$ matches $11.57$ to $0.6\%$. This provides a cleaner $E_8$-based structure for the half-integer exponents, though the specific pairing rules remain underived.

\textbf{Wolfenstein parametrization (EXP-191).} Setting $\lambda = \phi^{-3}$ in the standard Wolfenstein expansion, the remaining parameter $A$ also admits a geometric expression:
\begin{equation}
A = \phi^{-1/2} = \frac{1}{\sqrt{\phi}} = 0.7862 \quad \text{(exp: } 0.790 \pm 0.012\text{, error 0.5\%)}
\label{eq:wolfenstein_A}
\end{equation}
With both $\lambda$ and $A$ fixed geometrically, the Wolfenstein CKM depends on only two remaining parameters ($\bar{\rho}$, $\bar{\eta}$).

\textbf{Classification:} The Cabibbo angle with QCD correction is the strongest CKM prediction ($0.001\%$, 0 free parameters). The Wolfenstein $A = 1/\sqrt{\phi}$ is a new geometric result ($0.5\%$). The full CKM matrix is PATTERN+FIT ($3{-}6\%$ with 0--2 free parameters). A complete derivation from first principles is not yet achieved.

\textbf{CP-violating phase (EXP-170).} The CKM CP-violating phase $\delta_{\text{CKM}}$ admits a striking geometric expression:
\begin{equation}
\delta_{\text{CKM}} = \arctan\!\left(\sqrt{5}\right) = 65.91° \quad \text{(exp: } 65.5 \pm 2.5°\text{)}
\end{equation}
with error $0.55\%$ ($< 1\sigma$). The argument $\sqrt{5} = 2\phi - 1$ connects directly to the 600-cell spectrum: the difference between the largest and smallest eigenvalues is $\lambda_1 - \lambda_8 = 12 - (6 - 6\phi) = 6 + 6\phi = 6\sqrt{5}$. An alternative expression $\delta_{\text{CKM}} = 5\,\theta_C = 5 \times 13.28° = 66.41°$ gives $1.32\%$ error (also $< 1\sigma$), linking the CP phase to the Cabibbo angle through the graph diameter.

\textbf{Status: PATTERN} (algebraic $\sqrt{5}$ origin identified but not uniquely derived).

\textbf{Unitarity triangle angles (EXP-191).} The three angles of the CKM unitarity triangle admit phi-based expressions:
\begin{equation}
\alpha_{\text{UT}} = \arctan(\phi^5) = 84.85° \;\;(0.6°), \quad
\beta_{\text{UT}} = \arctan(\phi^{-2}) = 20.91° \;\;(1.3°), \quad
\gamma_{\text{UT}} = \frac{2\pi}{5} = 72.0° \;\;(0.4°)
\end{equation}
where the parenthetical values are deviations from experiment. The angle $\gamma = 2\pi/5 = 72°$ is the interior angle of a regular \textbf{pentagon}---the fundamental face of the icosahedron dual to the 600-cell---and lies well within $1\sigma$ of the measured value $72.4 \pm 5.0°$. The sum $\alpha + \beta + \gamma = 177.76°$ deviates from the unitarity constraint $180°$ by $1.2\%$, which is consistent with the individual approximations.

\textbf{Bare exponent origin (EXP-194).} The exponents $\{3, 7, 12\}$ can be expressed through the intersection numbers $a_1 = 5$ and $b_1 = 6$ of the 600-cell distance-regular graph:
\begin{equation}
n_{12} = 2a_1 - b_1 - 1 = 3, \qquad n_{23} = b_1 + 1 = 7, \qquad n_{13} = a_1 + b_1 + 1 = 12
\end{equation}
Equivalently, the spacings $\{4, 5\} = \{\dim(\mathbb{R}^4),\, a_1\}$ generate the sequence recursively: $3 \to 3+4=7 \to 7+5=12$. The first two exponents are even-index Lucas numbers: $3 = L_2$, $7 = L_4$. \textbf{Status:} STRONG PATTERN (the specific linear combinations are not yet derived from a variational principle).

\textbf{Quark--lepton complementarity (EXP-195).} The bare 600-cell angles satisfy an \emph{exact} algebraic identity:
\begin{equation}
\arctan(\phi^{-3}) + \arctan(\phi^{-1}) = \frac{\pi}{4}
\label{eq:qlc}
\end{equation}
since $\phi^{-3} + \phi^{-1} + \phi^{-4} = 1$, which is the condition for $\arctan(a) + \arctan(b) = \pi/4$. This means the bare Cabibbo angle plus the bare PMNS $\theta_{12}$ (Golden Ratio Mixing) sum to exactly $45°$---a geometric identity that holds independently of any fitting or approximation.

\subsection{PMNS Matrix (Neutrinos)}

\textbf{Geometric $\sin^2$ formulas (EXP-206--207).} The three PMNS mixing angles admit purely geometric expressions in terms of 600-cell invariants, \emph{without} free parameters:

\begin{table}[H]
\centering
\begin{tabular}{lccccl}
\toprule
\textbf{Angle} & \textbf{Formula} & \textbf{Predicted} & \textbf{PDG 2024} & \textbf{Error} & \textbf{Geometric origin} \\
\midrule
$\sin^2\theta_{13}$ & $\dfrac{1}{45} = \dfrac{1}{d \times N_{\text{eig}}}$ & 0.02222 & $0.02219 \pm 0.00062$ & \textbf{0.16\%} & $d=5$, $N_{\text{eig}}=9$ \\[8pt]
$\sin^2\theta_{12}$ & $\dfrac{2}{\phi + 5}$ & 0.3028 & $0.303 \pm 0.012$ & \textbf{0.33\%} & $2/(a_1 + \phi)$ \\[8pt]
$\sin^2\theta_{23}$ & $\dfrac{4}{7} = \dfrac{m(7)}{\Delta_{\text{Schur}}}$ & 0.5714 & $0.572 \pm 0.018$ & \textbf{0.32\%} & mult(7)/Schur gap \\
\bottomrule
\end{tabular}
\caption{PMNS mixing angles from 600-cell geometry. All three lie within $1\sigma$ of the PDG 2024 central values. $d$ = graph diameter, $N_{\text{eig}}$ = number of distinct eigenvalues, $a_1 = 5$ = first intersection number.}
\label{tab:pmns_geometric}
\end{table}

\textbf{Mass-squared ratio.} The ratio of neutrino mass-squared differences admits a coupling-constant expression:
\begin{equation}
\frac{\Delta m^2_{21}}{\Delta m^2_{31}} = 4\alpha \qquad \text{(error: 1.24\%)}
\end{equation}
With a radiative correction $\alpha_s/(12\pi)$ applied to $\theta_{12}$, the agreement improves to $0.01\%$.

\textbf{Phi-tower origin.} The angles also admit phi-tower representations with exponents $n \in \{0, 1, 4\}$:
\begin{itemize}
    \item $n=0$: $\phi^0 = 1$ (maximal mixing), $\theta_{23} \approx 45^\circ + \epsilon$
    \item $n=1$: $1/\phi$ (fundamental 600-cell ratio), $\theta_{12} = \arctan(1/\phi)$
    \item $n=4$: $\dim(\mathbb{R}^4)$, $\theta_{13} = \arctan(\phi^{-4})$
\end{itemize}

\textbf{Key insight:} Quarks ``see'' algebraic structure (quaternions, octonions), while leptons ``see'' geometric structure (space, ratios). The factor-3 rule $n^{\text{CKM}} = 3 \times n^{\text{PMNS}}$ holds exactly for $\theta_{12}$ and $\theta_{13}$.

\textbf{Status:} The three $\sin^2$ formulas are at the \textbf{PATTERN} level---all numerically within $1\sigma$, with geometrically motivated denominators, but not yet derived from a variational principle.

\subsection{CKM--PMNS Factor-3 Rule (EXP-171)}

A systematic comparison of CKM and PMNS exponents reveals a factor-of-3 relationship:
\begin{equation}
n^{\text{CKM}}_{ij} = 3 \times n^{\text{PMNS}}_{ij}
\end{equation}

\begin{center}
\begin{tabular}{lccc}
\toprule
\textbf{Angle} & $n^{\text{CKM}}$ & $n^{\text{PMNS}}$ & \textbf{Ratio} \\
\midrule
$\theta_{12}$ & 3 & 1 & \textbf{3} \\
$\theta_{13}$ & 12 & 4 & \textbf{3} \\
$\theta_{23}$ & 7 & 0 & $7 \neq 3 \times 0$ \\
\bottomrule
\end{tabular}
\end{center}

The factor 3 holds exactly for $\theta_{12}$ and $\theta_{13}$, with $\theta_{23}$ as the anomalous case ($7 \neq 0$). The factor 3 itself has a structural origin: each fermion vertex on the 600-cell graph has exactly 3 neighbors in each of the 4 other cosets (Section~\ref{sec:soliton_creation}), and $3 = \dim(\text{Im}(\mathbb{H})) = N_{\text{colors}}$.

\textbf{Ratio universality.} The ratio of exponents within each matrix is:
\begin{equation}
\frac{n_{13}}{n_{12}} = \frac{12}{3} = \frac{4}{1} = 4 \quad \text{(CKM and PMNS identical)}
\end{equation}
This inter-angle ratio of 4 is preserved across both sectors, suggesting a common geometric origin for the $\theta_{12}$--$\theta_{13}$ hierarchy.

\textbf{Status: PATTERN.} The factor 3 for $\theta_{12}$ and $\theta_{13}$ is exact, and the inter-angle ratio 4 is \textbf{DERIVED} (identical in both sectors). The anomalous $\theta_{23}$ and the physical origin of the factor 3 remain open.

%==============================================================
\begin{thebibliography}{99}

% === 600-Cell Theory ===

\bibitem{oneill2025gut}
C. C. O'Neill, ``The 600-Cell Grand Unified Theory: A Geometric and Group-Theoretic Model of the Standard Model and Beyond,'' preprint (2025).

\bibitem{oneill2025120cell}
C. C. O'Neill, ``The 120-Cell Theory of Everything,'' preprint (2025).

\bibitem{oneill2023binary}
C. C. O'Neill, ``Isomorphism of a Modified Standard Model Particle Assignment with the Binary Octahedral Group (24-cell),'' preprint (2023).

% === Geometric Optimality & Spectral Principles ===

\bibitem{cohn2007}
H. Cohn and A. Kumar, ``Universally optimal distribution of points on spheres,'' \textit{J. Amer. Math. Soc.} \textbf{20}, 99--148 (2007).

\bibitem{chamseddine1997}
A. H. Chamseddine and A. Connes, ``The Spectral Action Principle,'' \textit{Commun. Math. Phys.} \textbf{186}, 731--750 (1997).

\bibitem{cohn2011}
H. Cohn and J. Woo, ``Three-point bounds for energy minimization,'' \textit{J. Amer. Math. Soc.} (2011).

% === Experimental Data ===

\bibitem{codata2022}
CODATA, ``Recommended Values of the Fundamental Physical Constants,'' NIST (2022).

\bibitem{fermi2009}
Fermi LAT Collaboration, ``A limit on the variation of the speed of light arising from quantum gravity effects,'' \textit{Nature} \textbf{462}, 331--334 (2009).

\bibitem{pdg2024}
Particle Data Group, ``Review of Particle Physics,'' \textit{Phys. Rev. D} \textbf{110}, 030001 (2024).

% === Lattice & Number Theory ===

\bibitem{conway1998icosian}
J. H. Conway and N. J. A. Sloane, ``Sphere Packings, Lattices and Groups,'' 3rd ed., Springer (1998). Ch.~8: The icosian lattice and $E_8$.

% === Discrete Scale Invariance ===

\bibitem{sornette1998}
D. Sornette, ``Discrete-scale invariance and complex dimensions,'' \textit{Phys. Rep.} \textbf{297}, 239--270 (1998).


% === Quasicrystal & Projection ===

\bibitem{elser1987}
V. Elser and N. J. A. Sloane, ``A highly symmetric four-dimensional quasicrystal,'' \textit{J. Phys. A: Math. Gen.} \textbf{20}, 6161--6168 (1987).

% === E8 Universal Optimality ===

\bibitem{cohn2022}
H. Cohn, A. Kumar, S. D. Miller, D. Radchenko, and M. Viazovska, ``Universal optimality of the $E_8$ and Leech lattices and interpolation formulas,'' \textit{Ann. Math.} \textbf{196}, 983--1082 (2022).

\end{thebibliography}

\end{document}
